% !TEX root = ../main.tex
\chapter{구현과 실증 결과}
\label{ch:results}

이 장에서는 \ref{ch:methodology}장의 설계를 아비트럼 원(Arbitrum One)에서 하나의 사례 연구로 실제로 돌려 봐요. 체인은 하나이고, 관측 기간(점수를 매길 자료를 모은 기간)도 하나예요. 같은 기간 비교에는 모두 스마트 계약(블록체인 위에서 저절로 실행되는 프로그램)의 매칭 코호트(모든 점수가 똑같이 점수를 매기는 계약 묶음) 하나를 써요. 두 홀드아웃(나중 기간을 떼어 두고 맞혀 보는 시험) 기간 W0과 W1에는 스펜더(spender, 허락을 받아 대신 쓰는 쪽) 코호트와 트레이더(거래하는 쪽) 코호트를 써요. 크기에 따른 실행 시간을 볼 때는 매칭 코호트에서 일부만 뽑은 부분 표본을 써요. 이 장은 \ref{sec:amount-revision}절의 수정 분석(등록 뒤에 금액을 달러로 바꿔 다시 한 분석)을 보고해요. 수정 분석에서는 모든 금액을 달러(USD)로 매겨요. 그리고 코호트, 그래프(점과 화살표로 이은 그림), 대리 지표(대신 재는 값), 스펜더 라벨(나중에 채점에 쓰는 정답 값)에는 다섯 토큰 WETH, USDC, USDT0, ARB, WBTC만 남겨요. 이렇게 고치기로 정한 것은 같은 기간, W0, W1의 등록 분석(미리 등록한 그대로 한 분석)의 결과를 안 뒤였어요. 그래서 이 장의 어떤 결과도 등록 분석의 결과가 아니에요. 규칙~A는 이 연구의 로그(계약이 남기는 기록)를 하나도 뽑기 전에, 2026년 10월 9일의 분석 계획에서 정해 두었어요. 규칙~B와, 두 화살표(엣지) 종류 어느 쪽으로도 만들지 않은 결과 값에서의 비교는 나중에 등록했어요. 등록 분석의 W0 결과를 안 뒤였고, 이 연구가 W1의 로그를 뽑기 전이었어요. 이 장은 이 규칙들을 기준값과 문구까지 그대로, 고친 점수에 적용해요. 규칙~A는 W0과 같은 기간에, 규칙~B는 W1에 적용해요. 그 결과 값 비교는 W1에 적용하고, W0을 보조 기간으로 써요(\ref{sub:fresh-results}절과 \ref{sub:neutral-label-results}절). 여기서 나온 판정은 등록한 시험이 아니라 재분석이에요. 등록 분석도 같은 판정에 이르렀어요(부록~\ref{sec:map-registered}, 표~\ref{tab:registered-vs-revised}). 모든 표는 \ref{sec:reproducibility}절의 평가 요약 파일에서 만들어요. 켄달의 타우 $\tau$(두 줄 세우기가 얼마나 같은 순서인지 나타내는 수)에는 모두 95\% 신뢰구간(믿을 만한 범위)이 함께 나와요. 다만 금액 기울기 점검과 표본 정의 점검은 구간 없이 계산했어요. 그리고 페이지랭크(PageRank)는 모두 그것이 매끄럽게 다듬는(평활하는) 원시 차수(화살표를 그냥 센 수)와 견주어 읽어요. 각 절은 그 절의 비교 가운데 어느 것을 계획이나 등록 문서가 정해 두었는지 밝혀요.

\dissertationSubheading[sec:dataset]{사례 연구 배경: 아비트럼 원}

아비트럼 원을 고른 까닭은 세 가지예요. 첫째, 아비트럼 원은 이더리움 롤업(레이어 2)이고, 그 위에서 많은 디파이(DeFi, 탈중앙 금융) 서비스(프로토콜)가 돌아가요(\ref{sec:ethereum-l2}절). 둘째, 아비트럼 원의 로그는 공개된 빅쿼리(BigQuery) 내보내기 자료에 들어 있고, 처리 과정(파이프라인)이 바로 이 자료를 읽어요. 셋째, 아비트럼 원의 증거금 거래(빌린 돈을 보태 하는 거래) 서비스 가운데 하나인 GMX~V2가 포지션 닫기를 기록해요. 트레이더 라벨은 이 기록에서 나와요. 관측 기간은 2023년 10월 1일부터 2026년 6월 30일까지예요(33개월). 로지스틱 감쇠(오래된 일일수록 가볍게 세는 무게)는 이 기간의 끝인 $T_{\mathrm{obs}}$를 기준으로 해요. 기준 시점에 일어난 이벤트는 시간 무게가 $0.858$이에요. 무게는 180일에서 $0.5$로 떨어지고, 2년보다 오래된 이벤트에서는 0에 가까워요(그림~\ref{fig:logistic-decay-methods}). 앞쪽 몇 달의 이벤트도 화살표를 만들기는 해요. 하지만 AWP(적응형 가중 페이지랭크)와 엔도스랭크(EndorseRank)에서 화살표 무게의 대부분은 마지막 한 해가 차지해요.

\dissertationSubsubheading[sub:cohorts]{코호트}

계약 집합 네 개에 점수를 매기거나 평가해요. 이 집합에 든 주소는 모두 계약 코드를 가지고 있어요. 외부 소유 계정(EOA, 개인 키로 움직이는 계정)은 그래프에 주인(허락을 해 주는 쪽), 보낸 쪽(송신자), 받는 쪽(수신자)으로만 들어가고, 점수를 받지 않아요.

\begin{itemize}
\item 매칭 코호트($n=15{,}107$), \ref{sec:data-sources}절에서 정의했어요. 다섯 토큰 가운데 하나에 대해, 서로 다른 주인 세 명 이상에게서 0이 아닌 \texttt{Approval}을 받은 계약들이에요. 이 계약들은 계획의 코호트 27,844개의 일부예요. 계획의 코호트는 같은 규칙을 모든 토큰의 허락(allowance, 토큰을 얼마까지 대신 써도 된다는 허락)에 적용했어요. 그 코호트의 나머지 12,737개 계약은 다섯 토큰으로는 이 규칙을 채우지 못해요. 계획의 스펜더 후보 29,022개에 대해, 블록 513,216,506(2026년 10월 9일)에서 \texttt{eth\_getCode}를 불렀어요. 그러자 27,844개는 계약 코드를 돌려주었고, 1,121개는 코드가 없었고, 57개는 EIP-7702 위임 표시를 돌려주었어요\parencite{buterin2024eip7702}. 뒤의 두 무리는 EOA라서 뺐어요. 코호트에 드는지는 $T_{\mathrm{obs}}$의 어느 그래프에도 기대지 않아요. 코호트 계약 가운데 523개는 허락 그래프 밖에 있어요. 기간이 끝날 때, 그 계약들을 허락했던 모든 주인과 토큰의 마지막 \texttt{Approval} 값이 0이었기 때문이에요. 또 그 계약들은 주인으로서도 다섯 토큰의 양수 허락을 하나도 가지고 있지 않았어요. 송금(transfer) 그래프는 2,670개를 빼요. 기간 안에 다섯 토큰의 0이 아닌 송금 가운데, 이 계약들을 보낸 쪽이나 받는 쪽으로 적은 것이 하나도 없기 때문이에요(스펜더가 \texttt{transferFrom}으로 주인에게서 다른 세 번째 주소로 토큰을 바로 옮기면, 그 결과로 생긴 \texttt{Transfer} 이벤트에 그 스펜더는 나오지 않아요). 그리고 계약 80개는 두 그래프 모두의 밖에 있어요. 어떤 그래프에 없는 계약은 그 그래프를 쓰는 방법에서 0점을 받아요. 같은 기간의 일치도, 차이 검정, 벤치마크(속도 재기)는 모두 이 코호트에서 계산해요.
\item 두 홀드아웃 기간의 스펜더 코호트. W0에서는 점수를 $t_1$(2026년 3월 31일, 23:59:59~UTC)에 동결해요(고정해요). 라벨은 2026년 4월 1일부터 6월 30일까지 세요. 이 코호트는 뽑아 낸 허락 자료에서, $t_1$에 다섯 토큰 가운데 하나의 양수인 마지막 허락(가장 최근 허락)을 하나 이상 가진 모든 스펜더 계약이에요. 이런 스펜더 17,167개 가운데 17,050개는 계약 코드를 가지고 있고, 104개는 코드가 없고, 13개는 EIP-7702 위임 표시를 가지고 있어요. 이 EOA 117개는 뺐어요. W1에서는 점수를 $T_{\mathrm{obs}}$(2026년 6월 30일, 23:59:59~UTC)에 동결하고, 라벨은 2026년 7월 1일부터 9월 30일까지 세요. 같은 규칙으로 스펜더 계약 18,152개가 나와요. 스펜더 코호트는 매칭 코호트의 부분 집합이 아니에요. 뽑아 낸 허락에는 코호트 계약이 주인으로서 해 준 허락도 들어 있어요. 그래서 그 계약들이 허락해 준 계약도 스펜더로 들어와요. 그리고 스펜더는 양수 허락을 해 준 주인이 한 명만 있으면 돼요.
\item 트레이더 코호트. 계약 코드를 가지고 있고, 동결일 그래프의 점(노드)이고(어느 층이든), 라벨 기간 안에 포지션을 세 개 이상 닫은 GMX~V2 계정들이에요. W0에서는 88개이고, 그 가운데 40개는 동결일까지 허락을 받은 적이 있었어요. W1에서는 70개이고, 그 가운데 37개가 그랬어요. 계약 계정들은 관측 기간에 GMX~V2 포지션 133,685개를 닫았고, 그 가운데 6,473개는 청산(손실이 너무 커져 거래가 강제로 끝나는 일)으로 닫혔어요. 2025년에는 88,002번 닫았고, 2026년 상반기에는 9,809번 닫았어요. 그리고 2026년 7월부터 9월까지는 6,805번 닫았고, 그 가운데 182번은 청산이었어요. 그래서 트레이더 결과는 모두 적은 수의 계약에 기대고, 구간이 넓어요.
\item 매칭 코호트의 부분 표본. SHA256으로 늘 똑같이 뽑히도록 정한 부분 표본(시드 \texttt{contract-benchmark-v1})이에요. 계약 수는 $n=1{,}000$, $2{,}000$, $5{,}000$, $10{,}000$개이고, 그다음은 15,107개인 코호트 전체예요(\ref{sub:expanded-pool}절). 각 단계는 뽑힌 계약에 닿는, 뽑아 낸 화살표를 남겨요. 그래서 화살표 수는 $n$과 함께 늘어나요. $n=1{,}000$에서는 허락 화살표 424,888개와 송금 화살표 1,692,513개이고, 마지막에는 그래프 전체가 돼요. 이 단계들은 크기에 따른 실행 시간(\ref{sub:runtime-scaling}절)과 표본 정의에 대한 민감도(조건을 바꿔 결과가 얼마나 움직이는지 보기, \ref{sub:sample-size}절)를 보는 데 써요.
\end{itemize}

뽑아 낸 행(기록 한 줄)의 대부분은 다섯 토큰의 것이에요. 관측 기간에 계획의 코호트의 에고 네트워크(코호트 계약에 닿는 화살표만 모은 그래프)에 든 이벤트 가운데, 수정 분석은 허락 224,954,202개(57.7\%)와 송금 1,149,168,538개(74.7\%)를 남겨요. 다른 토큰의 허락 164,202,859개와 송금 364,468,917개는 빼요. 그리고 다섯 토큰의 것이라도 수정 코호트의 계약이 하나도 끼지 않은 허락 709,821개와 송금 24,766,612개도 빼요. 남긴 송금에는 모두 가격이 있어요. $T_{\mathrm{obs}}$에서 양수인 마지막 허락 12,590,613개 가운데 5,763,757개(45.8\%)는 무제한 허락이에요. 그리고 6,776,289개(53.8\%)는 C-PR(결합 페이지랭크, 두 그래프를 함께 걷는 점수)의 허락 층에서 쓰는 상한 $U$(이보다 큰 금액은 이 값으로 자름)에 닿아요. 표본으로 뽑은 날짜 38개에서, DefiLlama 가격과 체인 위에서 읽은 Chainlink 가격 피드(블록체인에 올라오는 가격 정보)는 서로 달랐어요. 그 차이의 중앙값은 토큰마다 0.01\%에서 0.08\%였고, 가장 큰 차이는 0.72\%였어요. 부록~\ref{sec:data-tokens}의 표~\ref{tab:usd-tokens}에 이 수치들을 토큰마다, 주소와 가격 피드와 함께 실었어요.

$T_{\mathrm{obs}}$에서 허락 그래프에는 점 6,221,665개와 화살표 11,210,114개가 있어요. 송금 그래프에는 점 7,884,602개와 화살표 25,001,216개가 있어요. $t_1$에서 두 쌍 표에는 허락 쌍 10,650,258개와 송금 쌍 23,680,692개가 들어 있어요. $T_{\mathrm{obs}}$의 화살표 수는 엔도스랭크나 AWP 무게가 양수인 쌍의 수예요. 점 수에는 각각 2,602개와 2,246개의 주소도 들어 있어요. 이 주소들이 낀 쌍은 모두 금액이 아주 작아서, 부동소수점 계산(컴퓨터가 소수를 다루는 방식)에서 금액 무게가 0이 돼요. 엔도스랭크와 AWP의 풀이 프로그램(solver)이 쓰는 그래프는 이 주소들을 빼요(표~\ref{tab:benchmark-runtime}). 그래서 같은 계약들 위에서 송금 그래프의 화살표 수는 허락 그래프의 2.23배예요. 허락 그래프의 화살표 수는 송금 그래프의 44.8\%예요. 주인--스펜더 쌍은 다섯 토큰 가운데 하나의 마지막 허락이 양수인 동안에만 허락 화살표가 돼요. 이와 달리 기간 안에 이 토큰들의 0이 아닌 송금을 주고받은 주소 쌍은 모두 송금 화살표가 돼요. 허락 그래프의 점 가운데 6,208,769개는 양수 허락을 하나 이상 가진 주인이에요. 그래서 이 그래프는 거의 모두가, 훨씬 적은 스펜더들을 가리키는 주인들로 이루어져 있어요. 두 그래프는 같은 추출, 곧 수정 코호트의 에고 네트워크에서 나와요. 그래서 점수들 사이의 차이는 화살표와 다시 출발(재시작)에서 비롯돼요.

\dissertationSubsubheading[sub:implementation]{구현}

모든 알고리즘은 페이지랭크를 하나의 파이썬(Python) 구현으로 계산해요. 감쇠 계수(화살표를 따라갈 확률)는 $d=0.85$이고, $L_1$ 노름(차이의 크기를 모두 더한 값)으로 잰 허용 오차는 $\varepsilon=10^{-8}$이고, 반복(한 바퀴) 상한은 $I_{\max}=300$이에요. 엔도스랭크와 AWP는 화살표 목록과 재시작 벡터가 달라요. C-PR은 두 화살표 목록을 점마다 섞어요(식~\eqref{eq:coupled-transition}). 평가 실행은 한 번마다 점수, 대리 지표별 통계, 부트스트랩(다시 뽑기) 구간을 요약 파일 하나에 써요. 벤치마크는 잰 시간을 따로 자기 파일에 써요. 등록 분석은 같은 기간 결과와 W0 결과를, 관측 기간을 뽑아 낸 뒤 W1을 등록하기 전에 계산했어요. W1 결과는 7월부터 9월까지의 로그를 뽑아 낸 뒤에 계산했어요. 수정 분석의 결과는 모두 2026년 10월 10일 11:36~UTC부터 계산했어요. 수정을 커밋(저장소에 바뀐 내용을 기록함)한 뒤였고, 세 기간의 등록 분석의 결과를 모두 안 뒤였어요(\ref{sec:amount-revision}절). $d=0.85$에서 모든 풀이는 반복 상한보다 훨씬 일찍 수렴했어요(값이 더는 거의 바뀌지 않게 되었어요). $T_{\mathrm{obs}}$에서 엔도스랭크는 반복 69번, AWP는 92번, $\lambda=0.5$인 C-PR은 79번, C-PR ($\lambda=0$)은 92번이 필요했어요. 가장 긴 풀이 하나는 S-PR(보증 씨앗 페이지랭크)의 씨앗을 쓴 송금 걷기였고, 95번이 필요했어요(\ref{sec:benchmark-results}절). 부록~\ref{ch:appendix-eval}에는 실행 목록, 저장소 버전, 설정 파일, 요약 파일, 그리고 벤치마크 시간을 잰 컴퓨터가 실려 있어요.

\dissertationSubheading[sec:benchmark-results]{계산 벤치마크(주장~C1)}

주장~C1(이 연구가 내세우는 것)은 같은 입력에서 각 점수를 계산하는 데 얼마나 드는지에 대한 것이에요. 표~\ref{tab:benchmark-runtime}은 \ref{sec:computational-benchmark}절의 시간 재는 절차를 따라, 알고리즘마다 몇 가지를 보고해요. 대상은 $T_{\mathrm{obs}}$에서 매칭 코호트의 전체 그래프예요. 보고하는 것은 다음과 같아요. 시간을 재지 않는 준비 실행을 한 번 한 뒤 시간 잰 실행을 다섯 번 해서 얻은 평균 푸는 시간과 그 표준편차(SD, 값이 퍼진 정도), 준비 실행 풀이에서 가장 많이 쓴 메모리, 반복 횟수, 그리고 풀이 프로그램이 쓰는 그래프의 점과 화살표 수예요. 각 그래프와 설정은 한 번씩만 시간을 쟀어요. 같은 측정이 표~\ref{tab:benchmark-scaling}의 코호트 전체 행과 표~\ref{tab:robustness-damping}의 $d=0.85$ 행에 다시 나와요. 그리고 부록~\ref{ch:appendix-eval}의 표~\ref{tab:benchmark-runs}에 시간 잰 실행이 모두 실려 있어요.

\IfFileExists{results/tables/benchmark-runtime.tex}{% Generated by experiments/scripts/export_latex.py -- do not edit by hand.
% Source: experiments/data/2-processed-tables-and-evaluations/revised-usd/benchmark/benchmark.json
\begin{table}[htbp]
\centering
\caption[매칭 코호트의 전체 그래프에서 잰 계산 벤치마크.]{매칭 코호트의 전체 그래프에서 잰 계산 벤치마크(속도 재기)예요(계약 $n=15{,}107$개. AWP 대 엔도스랭크의 실행 시간 비율은 2.2$\times$이고, 이때 화살표 수의 비율은 2.2$\times$예요). 실행 시간은 시간을 재지 않는 준비 실행을 한 번 한 뒤, 시간을 잰 풀이 프로그램(solver) 호출 5번에서 실제 걸린 시간(초)의 평균이에요. 최대 메모리는 그 준비 실행에서만 재요. 표준편차(SD)는 같은 5번의 실행에서 구했어요. 시간을 잰 호출 한 번은 화살표 목록을 점수로 바꿔요. 여기에는 점에 번호 붙이기, 듬성듬성한 행렬 만들기, 행 정규화(각 행의 합을 1로 맞추기), 재시작 벡터 만들기, 거듭제곱 반복이 들어가요. $|V|$와 $|E|$는 각 방법의 풀이 프로그램이 쓰는 그래프의 점 수와 서로 다른 화살표 수예요. C-PR은 허락 층과 송금 층을 합친 그래프 위를 걸어요. 그래서 C-PR의 $|E|$는 두 층의 서로 다른 쌍을 모두 세요. S-PR은 송금 층에서 돌아가요. S-PR의 실행 시간과 반복 수에는, 재시작 벡터를 만들어 주는 허락 층을 푸는 계산도 들어 있어요. 모든 실행은 표~\ref{tab:benchmark-runs}에 있어요.}
\label{tab:benchmark-runtime}
\fitwidth{%
\begin{tabular}{lrrrrrr}
\toprule
방법 & 실행 시간 (초) & 표준편차 (초) & 최대 메모리 (MB) & 반복 수 & $|V|$ & $|E|$ \\
\midrule
EndorseRank & 10.294 & 0.601 & 879.7 & 69 & 6{,}219{,}063 & 11{,}210{,}114 \\
AWP & 22.367 & 1.825 & 1{,}712.6 & 92 & 7{,}882{,}356 & 25{,}001{,}216 \\
\midrule
C-PR ($\lambda=0.5$) & 36.162 & 2.112 & 1{,}768.1 & 79 & 9{,}151{,}010 & 31{,}552{,}359 \\
S-PR & 34.422 & 3.833 & 1{,}235.8 & 158 & 7{,}884{,}602 & 25{,}018{,}892 \\
\bottomrule
\end{tabular}
}
\end{table}
}{}

엔도스랭크는 허락 그래프(점 6,219,063개, 화살표 11,210,114개)를 평균 10.29초에 풀고, AWP는 송금 그래프(점 7,882,356개, 화살표 25,001,216개)를 22.37초에 풀어요. 그래서 AWP는 엔도스랭크보다 2.17배 오래 걸려요. 이때 화살표 수의 비는 2.23이에요. 듬성듬성한(희소) 행렬을 짜 맞추는 일과 거듭제곱 반복(같은 계산을 되풀이하기)을 한 번 돌리는 일은 둘 다 $|E|$에 비례해요(부록~\ref{sec:pr-sparsity}). 그리고 시간의 비는 화살표 수의 비에 가깝고, 그보다 조금 낮아요. AWP가 반복을 더 많이 하는데도 그래요($T_{\mathrm{obs}}$의 평가 실행에서처럼 92번 대 69번). 가장 많이 쓴 메모리는 엔도스랭크가 879.7~MB, AWP가 1,712.6~MB예요. 시간 잰 실행 다섯 번의 표준편차는 0.60초와 1.83초예요. 이 속도는 $|E|$에서 나와요. 같은 계약들 위에서 허락 그래프의 화살표 수는 송금 그래프의 44.8\%이고, 두 점수는 같은 풀이 프로그램을 써요.

두 혼합 점수 행은 이 속도 결과가 어디까지 맞고 어디서 멈추는지 보여 줘요. C-PR은 두 화살표 집합을 합친 것(점 9,151,010개, 화살표 31,552,359개) 위를 걷고, 36.16초가 걸려요. AWP의 1.26배인 화살표 위에서 AWP보다 1.62배 오래 걸리는 거예요. S-PR은 송금 층(점 7,884,602개, 화살표 25,018,892개) 위에서 돌고, 34.42초가 걸려요. AWP보다 1.54배 오래 걸려요. S-PR의 시간과 반복 횟수에는 재시작 벡터를 대 주는 허락 층의 풀이도 들어 있어요. $T_{\mathrm{obs}}$의 평가 실행에서 반복 158번은 그 풀이의 63번과, 씨앗을 쓴 송금 걷기의 95번을 더한 거예요. 그래서 비용이 낮은 것은 허락 그래프만 걷는 점수에만 해당하고, 그 까닭은 그 그래프의 크기예요. 송금 층을 걷는 알고리즘은, 혼자 걷든 허락 정보와 합쳐 걷든, 그 층의 화살표만큼 값을 치러요.

이 결과에는 주의할 점이 두 가지 있어요. 이 시간은 미리 만들어 둔 화살표 목록에서 풀이 프로그램을 부르는 시간만 담아요. 처음부터 모두 새로 계산하려면 읽어 들이기(파싱), 해독(디코딩), 화살표 목록 만들기가 더해져요. 이 일들은 모든 알고리즘이 전부 또는 일부를 함께 해요. 그리고 더 빠른 방법이 꼭 더 좋은 방법은 아니에요. 그래서 실행 시간은 뒤에 나오는 일치도 결과와 홀드아웃 결과와 견주어 읽어야 해요.

\dissertationSubsubheading[sub:runtime-scaling]{매칭 코호트 부분 표본에서의 실행 시간}

표~\ref{tab:benchmark-scaling}는 \ref{sub:cohorts}절의 각 단계에서 엔도스랭크와 AWP의 실행 시간을 보고해요. 두 실행 시간의 비, 두 화살표 수의 비, 그리고 그 단계에서 두 풀이 그래프의 점과 화살표 수도 함께 보고해요. 단계가 바뀌면 크기만이 아니라 화살표 수의 비도 바뀌어요. 허락 그래프와 견주면 송금 그래프의 화살표 수는 $n=1{,}000$에서 3.98배예요(1,692,513개 대 424,888개). $n=2{,}000$에서는 2.49배, $n=5{,}000$에서는 2.60배, $n=10{,}000$에서는 2.57배, 코호트 전체에서는 2.23배예요. 풀이 시간이 $|E|$를 따른다면, 시간의 비도 이 화살표 수의 비를 따라야 해요. 다섯 단계에서 시간의 비는 2.89, 4.55, 3.60, 2.70, 2.17이에요. $n=10{,}000$부터는 시간의 비가 화살표 수의 비와 0.13 안쪽으로 가까워요. 코호트 전체에서는 화살표 수의 비보다 조금 아래에 있어요. $n=5{,}000$에서는 시간의 비가 화살표 수의 비보다 1.00 커요. 엔도스랭크의 풀이가 1초도 안 걸리는 가장 작은 두 단계에서는, 시간의 비가 화살표 수의 비의 0.73배와 1.83배예요. 그래서 시간의 비는 가장 큰 두 그래프에서는 화살표 수의 비를 따르지만, 더 작은 세 그래프에서는 따르지 않아요. 그림~\ref{fig:runtime-scaling}은 두 실행 시간을 단계 크기에 따라 그려요. 각 그래프의 화살표 수는 그 점 옆에 적혀 있어요. 단계들은 한 코호트 그래프의 부분 표본이에요. 그래서 이 단계들은 그래프 전체까지 화살표가 늘면서 실행 시간이 어떻게 느는지 보여 주지만, 그보다 큰 그래프는 보여 주지 않아요.

\IfFileExists{results/tables/benchmark-scaling.tex}{% Generated by experiments/scripts/export_latex.py -- do not edit by hand.
% Source: experiments/data/2-processed-tables-and-evaluations/revised-usd/benchmark/benchmark.json
\begin{table}[htbp]
\centering
\caption[평가 집합의 크기에 따른 실행 시간.]{평가 집합의 크기에 따라 엔도스랭크(ER)와 AWP의 실행 시간이 어떻게 바뀌는지 보여 줘요. 평가 집합은 매칭 코호트에서 SHA256으로 늘 같은 결과가 나오게 뽑은 부분 표본이에요(시드 \texttt{contract-benchmark-v1}). 가장 큰 것은 계약 $n=15{,}107$개인 코호트 전체예요. 단계마다 뽑힌 계약에 닿는, 추출한 화살표만 남겨요. 그래서 두 풀이 프로그램 그래프의 $|V|$와 $|E|$를 단계마다 적었어요. AWP/ER 시간과 AWP/ER $|E|$는 같은 행에 있는 두 실행 시간의 비율과 두 화살표 수의 비율이에요. 실행 시간은 시간을 재지 않는 준비 실행을 한 번 한 뒤, 시간 잰 풀이 프로그램 호출 5번에서 실제 걸린 시간을 초 단위로 평균한 값이에요. $^{a}$코호트 전체 행은 표~\ref{tab:benchmark-runtime}의 측정값을 다시 쓴 것이고, 시간을 다시 재지 않았어요.}
\label{tab:benchmark-scaling}
\fitwidth{%
\begin{tabular}{rrrrrrrrr}
\toprule
$n$ & ER (초) & AWP (초) & AWP/ER 시간 & AWP/ER $|E|$ & ER $|V|$ & ER $|E|$ & AWP $|V|$ & AWP $|E|$ \\
\midrule
1{,}000 & 0.580 & 1.677 & 2.9$\times$ & 4.0$\times$ & 404{,}563 & 424{,}888 & 1{,}124{,}621 & 1{,}692{,}513 \\
2{,}000 & 0.939 & 4.272 & 4.6$\times$ & 2.5$\times$ & 1{,}112{,}404 & 1{,}263{,}471 & 1{,}917{,}597 & 3{,}148{,}826 \\
5{,}000 & 2.496 & 8.995 & 3.6$\times$ & 2.6$\times$ & 2{,}104{,}606 & 2{,}988{,}450 & 3{,}539{,}528 & 7{,}769{,}268 \\
10{,}000 & 6.220 & 16.821 & 2.7$\times$ & 2.6$\times$ & 4{,}291{,}078 & 6{,}673{,}167 & 6{,}247{,}634 & 17{,}180{,}184 \\
15{,}107$^{a}$ & 10.294 & 22.367 & 2.2$\times$ & 2.2$\times$ & 6{,}219{,}063 & 11{,}210{,}114 & 7{,}882{,}356 & 25{,}001{,}216 \\
\bottomrule
\end{tabular}
}
\end{table}
}{}

\IfFileExists{\figpath{runtime-scaling.pdf}}{%
\begin{figure}[htbp]
\centering
\includegraphics[width=0.9\textwidth]{\figpath{runtime-scaling.pdf}}
\caption[{뽑은 계약 수에 따른 엔도스랭크와 AWP의 실행 시간.}]{단계에 든 계약 수에 따라, 엔도스랭크와 AWP의 풀이 프로그램을 한 번 부를 때 실제 걸린 시간을 그렸어요. 두 축은 모두 로그 눈금이에요. 단계는 매칭 코호트에서 SHA256으로 늘 똑같이 뽑히도록 정한 부분 표본이에요(표~\ref{tab:benchmark-scaling}). 각 단계는 그 계약들에 닿는, 뽑아 낸 화살표를 남겨요. 점 옆의 글자는 그 그래프의 화살표 수 $|E|$예요. 가장 오른쪽 점은 계약 15,107개인 코호트 전체예요. 오차 막대는 시간 잰 실행 다섯 번의 표준편차예요.}
\label{fig:runtime-scaling}
\end{figure}%
}{}

\dissertationSubheading[sec:same-window-alignment]{같은 기간 일치도(주장~C2)}

주장~C2는 각 점수가, 그 점수를 계산한 기간 안에서 무엇과 잘 맞는지 물어요. 표~\ref{tab:alignment-hybrid}은 모든 점수의 지표 가족(비슷한 지표 묶음) 평균 켄달의 $\tau$를 송금, 허락, 시빌(가짜 주소를 많이 만들어 속이는 공격) 안정성 가족마다 보고해요. 값마다 95\% 짝지은 부트스트랩 구간이 함께 나와요. 원시 차수에는 행이 없어요. 이 기간에서는 원시 차수가 대리 지표이기 때문이에요. 송금 입차수(송금을 보내온 서로 다른 주소 수)는 송금 가족에 들고, 받은 허락 수(허락 입차수)는 허락 가족에 들어요. 규칙~A의 보조 부분(\ref{sec:holdout-protocol}절)은 계획이 정해 둔 단 하나의 같은 기간 비교예요. 여기서는 등록 분석의 결과를 안 뒤에 이 부분을 고친 점수에 적용해요. 그래서 그 판정은 재분석이에요. 아래의 짝지은 대비(같은 다시 뽑은 표본에서 잰 두 값의 차이)는 계획에서 정해 두지 않았고, 같은 기간 결과와 함께 계산했어요.

% Generated by experiments/scripts/export_latex.py -- do not edit by hand.
% Source: experiments/data/2-processed-tables-and-evaluations/revised-usd/same-window/eval_summary.json
\begin{table}[htbp]
\centering
\caption[이 연구의 모든 점수에 대한 같은 기간 지표 가족 평균 켄달의 $\tau$.]{이 연구의 모든 점수에 대해, 같은 기간의 지표 가족 평균 켄달의 $\tau$와 부트스트랩 구간을 보여 줘요(매칭 코호트, 계약 $n=15{,}107$개). 점수는 엔도스랭크와 AWP, 다섯 가지 $\lambda$ 값의 결합 연산자 C-PR, 그리고 보증 씨앗 걷기 S-PR이에요. $\lambda=1$은 허락 층만 걷는 것이에요(달러로 나타낸 마지막 허락, 각각 상한을 둠). $\lambda=0$은 송금 층만 걷는 것이에요(달러 금액을 AWP의 로지스틱 감쇠로 줄임). 둘 다 고르게 다시 출발해요. 원시 차수는 송금 가족과 허락 가족의 대리 지표이기 때문에 따로 행이 없어요. 활동 재시작을 쓰는 엔도스랭크는 표~\ref{tab:endorserank-restarts}에 있어요. 신뢰구간(CI)은 계약을 다시 뽑아 구한 95\% 백분위 부트스트랩 구간이에요(짝지은 다시 뽑기 400번, 시드 42).}
\label{tab:alignment-hybrid}
\fitwidth{%
\begin{tabular}{lccc}
\toprule
방법 & 송금 & 허락 & 시빌 \\
\midrule
EndorseRank & 0.158 [0.147, 0.169] & 0.585 [0.578, 0.592] & 0.007 [-0.001, 0.015] \\
AWP & 0.572 [0.564, 0.580] & 0.193 [0.182, 0.203] & 0.395 [0.387, 0.403] \\
\midrule
C-PR ($\lambda=1$) & 0.162 [0.151, 0.173] & 0.597 [0.590, 0.604] & 0.011 [0.002, 0.019] \\
C-PR ($\lambda=0$) & 0.659 [0.653, 0.666] & 0.225 [0.214, 0.235] & 0.404 [0.396, 0.413] \\
C-PR ($\lambda=0.25$) & 0.385 [0.374, 0.395] & 0.457 [0.449, 0.465] & 0.155 [0.147, 0.164] \\
C-PR ($\lambda=0.5$) & 0.365 [0.355, 0.375] & 0.477 [0.469, 0.485] & 0.142 [0.134, 0.151] \\
C-PR ($\lambda=0.75$) & 0.344 [0.333, 0.354] & 0.492 [0.484, 0.500] & 0.129 [0.121, 0.138] \\
S-PR & 0.629 [0.621, 0.636] & 0.323 [0.313, 0.333] & 0.398 [0.390, 0.406] \\
\bottomrule
\end{tabular}
}
\end{table}


점수들은 자기가 걷는 화살표에 따라 갈려요. 엔도스랭크는 허락 가족에서 $0.585$ $[0.578,0.592]$, 송금 가족에서 $0.158$ $[0.147,0.169]$에 이르러요. AWP는 송금 가족에서 $0.572$ $[0.564,0.580]$, 허락 가족에서 $0.193$ $[0.182,0.203]$에 이르러요. 두 점수는 저마다 자기 화살표의 가족에서 상대보다 앞서요. 엔도스랭크는 허락 가족에서 $+0.392$ $[0.380,0.405]$만큼, AWP는 송금 가족에서 $+0.414$ $[0.401,0.427]$만큼 앞서요(부록~\ref{ch:appendix-er-awp}의 표~\ref{tab:tau-diff}). C-PR의 두 층을 각각 한 층만 걸어도 같은 갈림이 나와요. C-PR ($\lambda=1$)은 허락 가족에서 $0.597$ $[0.590,0.604]$에 이르고, C-PR ($\lambda=0$)은 송금 가족에서 $0.659$ $[0.653,0.666]$에 이르러요. 시빌 안정성 가족에서는 송금 위를 걷는 점수들이 앞서요. AWP는 $0.395$ $[0.387,0.403]$이고, 엔도스랭크는 $0.007$ $[-0.001,0.015]$예요. 이 일치의 일부는 만든 방식 때문에 저절로 생겨요. 각 가족의 주된 대리 지표는 그래프 하나와 같은 이벤트에서 나와요. 송금 입차수와 달러로 잰 받은 양의 합은 다섯 토큰의 송금에서 나와요. 받은 허락 수와 상한으로 자른 금액은 다섯 토큰의 마지막 허락에서 나와요. 그리고 페이지랭크는 자기 그래프의 입차수를 매끄럽게 다듬어요. 그래서 점수가 자기 화살표 종류의 대리 지표와 맞는 것은 의도한 구성개념(점수가 재려고 하는 생각) 점검이에요. 이것을 외부 시험(바깥 자료로 하는 시험)으로 칠 수는 없어요. \ref{sec:holdout-results}절의 홀드아웃은 같은 질문을 기간 밖에서 물어요.

부록~\ref{sec:proxy-family-detail}는 엔도스랭크와 AWP의 지표 가족 평균을 대리 지표 하나하나로 나누어 보여 줘요. AWP는 들어온 송금의 달러 금액 합($0.566$ $[0.558,0.575]$)을, 보낸 쪽의 수인 입차수($0.578$ $[0.569,0.587]$)와 거의 같은 만큼 가깝게 따라가요. AWP의 무게를 보면 예상할 수 있는 결과예요. 기울기를 송금의 중앙값에 맞추었기 때문에(\ref{sub:awp-graph}절), 200달러 송금의 무게는 2분의 1이고, 400달러 송금은 $0.8$, 100달러 송금은 $0.27$이에요. 그래서 화살표 무게는 송금 횟수뿐 아니라 보낸 금액에 따라서도 커져요. 등록 분석에서는 원래 단위(소수점 조정을 하지 않은 가장 작은 단위)에서 $b=1$이었어요. 그래서 금액과 상관없이 거의 모든 송금에서 $V$가 1에 가까웠어요. 그리고 AWP는 받은 양의 합($0.353$ $[0.345,0.361]$)보다 입차수($0.541$ $[0.534,0.547]$)를 훨씬 가깝게 따라갔어요(등록 분석의 결과 요약, 부록~\ref{sec:map-registered}). 시빌 안정성 가족에서 AWP는 송금 활동 기간($0.508$ $[0.499,0.517]$)과 활동한 달 수($0.549$ $[0.540,0.558]$)를 따라가요. 그리고 들어오는 상대 비율에서는 양수예요($0.127$ $[0.114,0.141]$).

엔도스랭크는 받은 허락 금액($0.515$ $[0.506,0.524]$)보다 받은 허락 수($0.655$ $[0.648,0.662]$)를 더 가깝게 따라가요. 엔도스랭크의 금액 무게는 허락 금액의 중앙값인 62달러의 두 배에서 $0.8$이고, 몇백 달러에서는 1에 가까워요. 그래서 그런 허락과 무제한 허락은 무게가 거의 같아요. 이와 달리 받은 허락 금액은 허락 하나하나를 상한인 20,000달러까지 세요. 엔도스랭크의 송금 계수는 입차수와 $0.182$ $[0.169,0.194]$이고, 받은 양의 합과 $0.134$ $[0.123,0.145]$예요. 시빌 안정성 가족에서 엔도스랭크는 들어오는 상대 비율에서 음수이고($-0.156$ $[-0.166,-0.147]$), 활동 기간에서 $0.080$ $[0.069,0.091]$, 활동한 달 수에서 $0.098$ $[0.087,0.109]$에 이르러요. 그래서 한 계약을 몇 명의 주인이 허락하는지는, 그 계약이 얼마나 오랫동안 송금을 받아 왔는지에 대해 알려 주는 것이 적어요. C-PR ($\lambda=1$)은 상한으로 자른 달러 금액으로 같은 허락 그래프를 걷고, 마찬가지로 고르게 다시 출발해요(균등 재시작). 이 점수는 모든 가족에서 엔도스랭크와 가까워요. 받은 허락 금액은 C-PR의 허락 층과 똑같이 상한으로 자른 금액을 더해요. 이 지표에서 이 점수의 점 추정값(구간이 아니라 값 하나로 낸 추정)은 더 높아요($0.540$ $[0.531,0.549]$ 대 $0.515$ $[0.506,0.524]$, 짝지은 대비는 계산하지 않았어요).

코호트의 모든 계약은 다섯 토큰의 허락을 주인 세 명 이상에게서 받았어요. 그래서 여기서 엔도스랭크는 동점(점수가 똑같은 것)을 거의 남기지 않아요. 엔도스랭크는 계약 15,107개에 서로 다른 점수 13,625개를 줘요. 계약 쌍 가운데 0.17\%만 동점이고, 그 대부분은 0점을 받은 계약 535개 사이에 있어요(\ref{sec:rank-divergence}절). 그래서 엔도스랭크의 계수는 거의 모든 계약의 순서에 바탕을 둬요. 그리고 그래프 밖의 계약을 외톨이 점(화살표가 하나도 없는 점)으로 남겨 두어도 계수는 하나도 바뀌지 않아요(표~\ref{tab:tie-sensitivity}).

주장~C2는 의도한 구성개념 점검으로서 뒷받침돼요. 다만 그 일치의 일부는 만든 방식 때문에 저절로 생겨요. 모든 점수는 자기가 걷는 화살표 종류의 대리 지표와 가장 잘 맞고, 그 구간은 0에서 멀리 떨어져 있어요. 그리고 C-PR의 섞은 점수들은 모든 가족에서 두 층 사이에 있어요(\ref{sec:coupled-results}절).

\dissertationSubsubheading[sec:coupled-results]{결합 연산자(RQ5)}

\ref{sub:coupled-graph}절에서 혼합 점수 두 개를 정의했어요. 결합 연산자 C-PR은 각 점의 이동 행렬(전이 행렬) 줄에서, 달러로 잰 허락 층과 송금 층을 무게 $\lambda$로 섞어요. 보증 씨앗 걷기 S-PR은 송금 층을 그대로 쓰지만, 허락 걷기를 바탕으로 다시 출발해요. 표~\ref{tab:tau-diff-hybrid}는 이 기간에서 두 혼합 점수의 짝지은 대비를 보여 줘요. 먼저 여덟 개는 한 층만 걷는 두 점수와 견주고, 그다음 네 개는 엔도스랭크와 AWP와 견줘요. 연구 질문 RQ5의 답은 기간 밖에서 정해요(\ref{sec:holdout-results}절과 \ref{sub:fresh-results}절). 여기서는 한 층 점수들이 통과하도록 만들어진 점검 위에서, 혼합 점수들이 어디에 있는지만 살펴봐요.

% Generated by experiments/scripts/export_latex.py -- do not edit by hand.
% Source: experiments/data/2-processed-tables-and-evaluations/revised-usd/same-window/eval_summary.json
\begin{table}[htbp]
\centering
\caption[C-PR과 S-PR을 두 층과, C-PR을 엔도스랭크 및 AWP와 견준 같은 기간 짝지은 대비.]{같은 기간에서 구한 짝지은 대비예요. 결합 연산자(C-PR, $\lambda=0.5$)와 보증 씨앗 걷기(S-PR)를, 이 둘을 만드는 두 층과 견주었어요. 두 층은 허락 층(C-PR ($\lambda=1$))과 송금 층(C-PR ($\lambda=0$))이에요. 또 C-PR을 엔도스랭크 및 AWP와 견주었어요(매칭 코호트, 계약 $n=15{,}107$개). 값은 지표 가족 평균 $\tau$이고, $a$와 $b$에는 같은 계약 다시 뽑은 표본을 써요. 층과의 대비는 규칙 A의 보조 부분 옆에 두는 더 엄격한 점검이에요. 규칙 A의 보조 부분은 C-PR의 점 추정값을 표~\ref{tab:alignment-hybrid}에서 더 나은 층의 구간과 비교해요. 규칙 A의 보조 부분은 분석 계획에 미리 정해 두었어요. 여기서는 그 규칙을 수정 분석의 점수에 썼어요. 그 결과 보조 부분의 조건은 채워지지 않아요. C-PR의 점 추정값은 송금 가족에서 C-PR ($\lambda=0$)의 구간 밖에 있고, 허락 가족에서 C-PR ($\lambda=1$)의 구간 밖에 있어요. 그리고 시빌 안정성 가족에서 C-PR은 두 층을 모두 넘지는 못해요. 구간이 0보다 아래이면, 표본을 뽑을 때 생기는 흔들림을 넘어서 $a$가 $b$보다 낮다는 뜻이에요. 구간이 0보다 위이면, 그 흔들림을 넘어서 더 높다는 뜻이에요. 신뢰구간(CI)은 계약을 다시 뽑아 구한 95\% 백분위 부트스트랩 구간이에요(짝지은 다시 뽑기 400번, 시드 42).}
\label{tab:tau-diff-hybrid}
\fitwidth{%
\begin{tabular}{lccccc}
\toprule
대비 ($a$ 빼기 $b$) & $\tau_a$ & $\tau_b$ & $\Delta\tau$ & $\Delta\tau$의 95\% 신뢰구간 & $P(\Delta\tau>0)$ \\
\midrule
\multicolumn{6}{l}{\emph{한 층만 걷는 두 층과 비교}} \\
C-PR 송금 빼기 C-PR ($\lambda=0$) 송금 & 0.365 & 0.659 & -0.294 & [-0.305, -0.284] & 0.000 \\
C-PR 허락 빼기 C-PR ($\lambda=1$) 허락 & 0.477 & 0.597 & -0.120 & [-0.127, -0.113] & 0.000 \\
C-PR 시빌 안정성 빼기 C-PR ($\lambda=0$) 시빌 안정성 & 0.142 & 0.404 & -0.262 & [-0.273, -0.253] & 0.000 \\
C-PR 시빌 안정성 빼기 C-PR ($\lambda=1$) 시빌 안정성 & 0.142 & 0.011 & +0.131 & [0.126, 0.137] & 1.000 \\
S-PR 송금 빼기 C-PR ($\lambda=0$) 송금 & 0.629 & 0.659 & -0.031 & [-0.035, -0.028] & 0.000 \\
S-PR 허락 빼기 C-PR ($\lambda=1$) 허락 & 0.323 & 0.597 & -0.274 & [-0.285, -0.263] & 0.000 \\
S-PR 시빌 안정성 빼기 C-PR ($\lambda=0$) 시빌 안정성 & 0.398 & 0.404 & -0.007 & [-0.009, -0.004] & 0.000 \\
S-PR 시빌 안정성 빼기 C-PR ($\lambda=1$) 시빌 안정성 & 0.398 & 0.011 & +0.387 & [0.376, 0.399] & 1.000 \\
\midrule
\multicolumn{6}{l}{\emph{EndorseRank, AWP와 비교}} \\
C-PR 송금 빼기 AWP 송금 & 0.365 & 0.572 & -0.207 & [-0.220, -0.194] & 0.000 \\
C-PR 허락 빼기 EndorseRank 허락 & 0.477 & 0.585 & -0.108 & [-0.115, -0.101] & 0.000 \\
C-PR 시빌 안정성 빼기 AWP 시빌 안정성 & 0.142 & 0.395 & -0.253 & [-0.263, -0.242] & 0.000 \\
C-PR 시빌 안정성 빼기 EndorseRank 시빌 안정성 & 0.142 & 0.007 & +0.135 & [0.129, 0.141] & 1.000 \\
\bottomrule
\end{tabular}
}
\end{table}


C-PR은 모든 가족에서 자기의 두 층 사이에 있어요. 섞는 비율 $\lambda=0.5$에서 C-PR은 송금 가족에서 $0.365$ $[0.355,0.375]$, 허락 가족에서 $0.477$ $[0.469,0.485]$, 시빌 안정성 가족에서 $0.142$ $[0.134,0.151]$에 이르러요. 그리고 $\lambda$가 $0.25$에서 $0.75$로 커지면, 세 가족 모두에서 허락 층 쪽으로 움직여요(표~\ref{tab:alignment-hybrid}). C-PR은 송금 층보다 허락 층에 더 가까이 있어요. 허락 가족에서는 C-PR ($\lambda=1$)보다 $0.120$ 낮고($[-0.127,-0.113]$), 송금 가족에서는 C-PR ($\lambda=0$)보다 $0.294$ 낮아요($[-0.305,-0.284]$). 식~\eqref{eq:coupled-transition}의 섞기 규칙은 두 층 모두에 나가는 화살표가 있는 점마다 작동해요. 허락 층의 점은 거의 모두 나가는 화살표가 있는 주인이에요. 그래서 주인이 송금도 보냈고 추출에서 그 송금을 남겼다면, 그곳마다 허락 줄이 걷기에 들어가요. 한 층 점수들과 견주면, C-PR은 송금 가족에서 AWP보다 낮고($-0.207$, $[-0.220,-0.194]$), 허락 가족에서 엔도스랭크보다 낮아요($-0.108$, $[-0.115,-0.101]$). 시빌 안정성 가족에서는 허락 걷기($+0.131$, $[0.126,0.137]$)와 엔도스랭크($+0.135$, $[0.129,0.141]$)보다 높고, 송금 걷기($-0.262$, $[-0.273,-0.253]$)와 AWP($-0.253$, $[-0.263,-0.242]$)보다 낮아요.

S-PR은 모든 가족에서, 그 가족에서 앞선 한 층 점수보다 낮아요. 송금 걷기를 허락 걷기를 바탕으로 다시 출발하게 하면, 송금 계수는 $0.629$ $[0.621,0.636]$로 낮아져요(C-PR ($\lambda=0$)과 견주어 $\Delta\tau=-0.031$, $[-0.035,-0.028]$). 시빌 안정성 계수도 $0.398$ $[0.390,0.406]$로 낮아져요($-0.007$, $[-0.009,-0.004]$). 허락 계수는 $0.323$ $[0.313,0.333]$로 올라가요. 고르게 다시 출발하는 송금 걷기는 $0.225$ $[0.214,0.235]$예요(이 둘의 짝지은 대비는 계산하지 않았어요). 하지만 허락 계수는 여전히 허락 걷기보다 $0.274$ 낮아요($[-0.285,-0.263]$). 재시작 벡터는 순간 이동(다른 곳으로 훌쩍 건너뛰기)을 할 때마다 걷기가 어디에 내려앉는지를 바꿔요. 이 벡터는 허락 신호의 일부를 송금 걷기로 옮겨 오지만, 허락 화살표를 직접 걷는 것보다는 적게 옮겨 와요.

분석 계획에서 정한 대로, 규칙~A의 보조 부분은 다음을 요구해요. 송금 가족과 허락 가족에서는 $\lambda=0.5$인 C-PR의 점 추정값이 더 나은 층의 95\% 구간 안에 있어야 해요. 그리고 시빌 안정성 가족에서는 C-PR이 두 층을 모두 넘어야 해요. 고친 점수에 적용하면, 세 부분이 모두 실패해요. 등록 분석에서도 그랬어요(표~\ref{tab:registered-vs-revised}). 송금 계수($0.365$)는 C-PR ($\lambda=0$)의 구간($[0.653,0.666]$)보다 아래에 있어요. 허락 계수($0.477$)는 C-PR ($\lambda=1$)의 구간($[0.590,0.604]$)보다 아래에 있어요. 그리고 시빌 안정성 가족에서 C-PR은 허락 층만 넘어요. 표~\ref{tab:tau-diff-hybrid}의 짝지은 대비도 같은 쪽을 가리켜요. 이 대비는 규칙에 들지 않는, 더 엄격한 점검이에요. 같은 기간에서 섞은 점수는 각 층이 자기 가족과 맞던 정도의 일부를 잃고, 어느 가족에서도 더 나은 층을 넘지 못해요. 남은 질문은 홀드아웃이 맡아요. 섞은 점수에 한 층 점수들에는 없는 예측(미리 맞히는) 정보가 있는지 하는 질문이에요.

\dissertationSubheading[sec:rank-divergence]{순위 갈라짐(H1 / RQ1)}

가설~H1과 연구 질문 RQ1은 엔도스랭크와 AWP가 같은 계약들을 서로 다른 순서로 세우는지를 물어요. 두 점수 벡터 사이의 방법 간 켄달의 $\tau_b$는 $0.174$이고, 구간은 $[0.163,0.185]$예요(표~\ref{tab:tie-sensitivity}의 마지막 행). 그래서 이 구간에는 0도 1도 들어 있지 않아요. 두 순위는 서로 관계가 있지만, 서로 바꿔 쓸 수 있는 것과는 거리가 멀어요. 계약 쌍 114,103,171개 가운데 96.7\%는 두 점수 모두에서 순서가 정해져 있어요. 그리고 그 가운데 41.2\%는 두 점수에서 순서가 반대예요.

이 결과에서 동점은 거의 아무 역할도 하지 않아요. 각 점수를 하나씩 따로 본 분포(주변 분포)는 그림~\ref{fig:score-distributions}에 있어요. 엔도스랭크는 계약 15,107개에 서로 다른 점수 13,625개를 주고, 쌍의 0.17\%를 동점으로 둬요. AWP는 서로 다른 점수 12,435개를 주고, 3.13\%를 동점으로 둬요. 두 점수 모두에서 가장 큰 동점 덩어리는 0점을 받은 계약들이에요. 엔도스랭크에서는 535개, AWP에서는 2,672개예요. 이 계약들은 각 그래프 밖에 있는 523개와 2,670개(그 가운데 80개는 두 그래프 모두의 밖), 그리고 금액이 너무 작아 금액 무게가 0과 다르지 않은 쌍만 가진 12개와 2개예요. 엔도스랭크에서는 계약 166개가 더 하나의 점수 $c=7.64\times10^{-8}$을 함께 가져요. 이 점수는 들어오는 화살표가 없는 점의 점수예요. 이 계약들은 허락 그래프에 주인으로만 들어 있어요. 나머지 코호트 계약은 모두 양수 허락을 하나 이상 받았고, $c$보다 높은 점수를 받아요. 주변 분포로는 두 점수가 같은 계약들을 어떻게 줄 세우는지 보여 줄 수 없어요. 그래서 그림~\ref{fig:rank-scatter}은 무작위로 뽑은 계약 3,000개에서 두 순위를 서로 맞대어 그려요. 동점인 계약은 $\tau_b$에서처럼 중간 순위(동점끼리 나눠 가진 순위의 평균)에 놓아요. 0점 덩어리는 오른쪽 끝의 세로줄(엔도스랭크)과 맨 위의 가로줄(AWP)로 나타나요. 두 점수가 모두 위쪽에 두거나 모두 아래쪽에 두는 계약은 항등선(두 순위가 같은 곳을 이은 대각선) 가까이에 있어요. 그 사이에서는 점들이 사각형의 대부분에 흩어져 있어요.

0점 덩어리는 처리 과정의 규칙 하나에서 나와요. 그래프에 없는 계약은 0점을 받는다는 규칙이에요(\ref{sec:graph-construction}절). 분석 계획은 다른 방법, 곧 없는 계약을 외톨이 점으로 남겨 두는 방법을 민감도 분석으로 정해 두었어요. 표~\ref{tab:tie-sensitivity}는 고친 점수에서 그 결과를 보여 줘요. 소수점 셋째 자리에서 바뀌는 계수는 하나도 없어요. 그리고 두 점수 사이의 일치는 $0.174$ $[0.163,0.185]$ 그대로예요. 그 까닭은 없는 계약들이 어디에 놓이는지에 있어요. 엔도스랭크는 고르게 다시 출발하기 때문에, 외톨이 점은 그래프의 점이 가질 수 있는 가장 낮은 점수인 $c$를 받아요. 그래서 허락 그래프 밖의 계약 523개는 순서의 맨 아래 가까이에 머물고, $c$를 받은 계약 166개와 한데 모여요. 무게가 0인 쌍만 가진 계약 12개는 허락 그래프의 점이에요. 그래서 이 규칙은 그 계약들을 더하지 않아요. 그리고 무게가 양수인 화살표만 남기는 풀이 프로그램은 그 계약들을 빼요. 그래서 그 계약들은 0점 그대로예요. 방향이 바뀌는 쌍은 하나도 없어요. 523개와 166개 사이의 쌍은 동점이 돼요. 그리고 523개와 12개 사이의 쌍은 주된 규칙에서는 0점으로 동점이었는데, 이제는 순서가 생겨요. AWP의 다시 출발 규칙에서는, 아무것도 보내지 않은 외톨이 계약은 다시 출발 몫을 받지 못해요. 그래서 두 규칙 모두에서 0점이에요. 코호트는 받은 허락으로 정의했어요. 그래서 거의 모든 코호트 계약은 허락 그래프에서 스펜더이고, 없는 계약을 다루는 규칙은 거의 아무것도 정하지 않아요.

% Generated by experiments/scripts/export_latex.py -- do not edit by hand.
% Source: experiments/data/2-processed-tables-and-evaluations/revised-usd/same-window/eval_summary.json
\begin{table}[htbp]
\centering
\caption[빠진 계약을 0점으로 매기거나 외톨이 점으로 둘 때의 같은 기간 지표 가족 평균.]{그래프에서 빠진 계약을 0점으로 매기거나 외톨이 점(화살표가 하나도 없는 점)으로 둘 때, 같은 기간의 지표 가족 평균 켄달의 $\tau$를 보여 줘요(매칭 코호트, 계약 $n=15{,}107$개). 주 분석에서는 어떤 방법의 그래프에 없는 계약을 0점으로 매겨요. 그 계약을 외톨이 점으로 두는 것은 분석 계획에 미리 정해 둔 민감도 분석이에요. 이 규칙의 영향을 받는 계약은 엔도스랭크에서 523개, AWP에서 2{,}670개예요(80개는 두 그래프 모두에 없어요). 엔도스랭크는 고르게 다시 출발하기 때문에, 외톨이 점은 들어오는 화살표가 없는 점과 같은 점수를 받아요. AWP는 송금을 보낸 주소에서만 다시 출발해요. 그래서 외톨이 계약은 점수 0을 그대로 갖고, AWP의 계수는 두 규칙에서 똑같아요. 마지막 행은 두 점수 사이의 일치예요. 신뢰구간(CI)은 계약을 다시 뽑아 구한 95\% 백분위 부트스트랩 구간이에요(짝지은 다시 뽑기 400번, 시드 42).}
\label{tab:tie-sensitivity}
\fitwidth{%
\begin{tabular}{lcccc}
\toprule
 & \multicolumn{2}{c}{빠진 계약을 0점으로} & \multicolumn{2}{c}{빠진 계약을 외톨이 점으로} \\
\cmidrule(lr){2-3} \cmidrule(lr){4-5}
지표 가족 & EndorseRank & AWP & EndorseRank & AWP \\
\midrule
송금 & 0.158 [0.147, 0.169] & 0.572 [0.564, 0.580] & 0.158 [0.146, 0.169] & 0.572 [0.564, 0.580] \\
허락 & 0.585 [0.578, 0.592] & 0.193 [0.182, 0.203] & 0.585 [0.579, 0.592] & 0.193 [0.182, 0.203] \\
시빌 안정성 & 0.007 [-0.001, 0.015] & 0.395 [0.387, 0.403] & 0.007 [-0.002, 0.015] & 0.395 [0.387, 0.403] \\
\midrule
EndorseRank 대 AWP & \multicolumn{2}{c}{0.174 [0.163, 0.185]} & \multicolumn{2}{c}{0.174 [0.163, 0.185]} \\
\bottomrule
\end{tabular}
}
\end{table}


\begin{figure}[htbp]
\centering
\includegraphics[width=0.9\textwidth]{\figpath{score-distributions.pdf}}
\caption[{매칭 코호트(계약 $n=15{,}107$개)에서 엔도스랭크와 AWP 점수의 주변 분포(점수가 양수인 계약, $\log_{10}$ 눈금).}]{매칭 코호트(계약 $n=15{,}107$개)에서 엔도스랭크와 AWP 점수의 주변 분포를 그렸어요. 점수가 양수인 계약을 $\log_{10}$ 눈금으로 그렸고, 개수도 로그 눈금이에요. 0점인 계약은 뺐어요. 이 계약들에는 그 그래프에서 무게가 양수인 화살표가 하나도 닿지 않아요. 엔도스랭크에서는 $535$개, AWP에서는 $2{,}672$개예요. 엔도스랭크에서 허락 그래프에 주인으로만 들어 있는 계약 $166$개는, 들어오는 화살표가 없는 점의 점수 $c=7.64\times10^{-8}$을 함께 가져요. 이것이 가장 낮은 양수 점수예요.}
\label{fig:score-distributions}
\end{figure}

\begin{figure}[htbp]
\centering
\includegraphics[width=0.8\textwidth]{\figpath{rank-scatter.pdf}}
\caption[{매칭 코호트에서 엔도스랭크와 AWP의 순위--순위 산점도(무작위로 뽑은 계약 $3{,}000$개).}]{매칭 코호트에서 엔도스랭크와 AWP의 순위를 서로 맞대어 찍은 산점도(점을 흩어 찍은 그림)예요(계약 $15{,}107$개 가운데 무작위로 뽑은 $3{,}000$개, 시드 42). 점수가 같으면 순위의 평균을 함께 가져요. 이것이 켄달의 $\tau_b$가 동점을 다루는 방법이에요. 오른쪽 끝의 세로줄에는, 엔도스랭크에서 0점인 $535$개 가운데 뽑힌 계약들이 있어요. 맨 위의 가로줄에는, AWP에서 0점인 $2{,}672$개 가운데 뽑힌 계약들이 있어요. 점선으로 그린 항등선 위의 점은 두 순서에서 순위가 같아요. 범례(그림 속 설명 상자)에는 모든 계약에서 구한 $\tau_b$와 그 95\% 구간이 있어요.}
\label{fig:rank-scatter}
\end{figure}

H1은 뒷받침돼요. 두 점수는 코호트를 서로 다르게 줄 세워요. 그리고 그 차이는 동점에서도, 없는 계약을 다루는 규칙에서도 나오지 않아요. 송금에 바탕을 둔 점수에 기대는 서비스는, 허락에 바탕을 둔 점수로 바꾸거나 그것과 섞어도 계약의 순서가 그대로 남으리라고 가정할 수 없어요.

\dissertationSubheading[sec:robustness-results]{강건성 점검과 민감도 점검}

엔도스랭크와 AWP에 대해 세 가지 점검을 했어요. 세 점검은 감쇠 계수, 금액 무게의 기울기, 표본을 정하는 방법을 각각 바꾸어 봐요. 감쇠 계수 값들은 분석 계획에서 미리 정했어요. 금액 기울기 점검은 계획에 있던 토큰 점검을 대신해요. 모든 그래프가 다섯 토큰을 남기게 되자, 토큰 점검은 할 까닭이 없어졌기 때문이에요(\ref{sub:top-tokens}절). 표본 정의 점검은 계획에 없었어요. 세 점검은 모두 등록 분석의 결과를 안 뒤에 했고, 고친 점수로 했어요. 처음 두 점검은 표본을 그대로 두고, 결과가 어떤 한 설정에 달려 있는지 물어요. 세 번째 점검은 표본 자체를 바꿔요. 그래서 계수가 움직일 것으로 예상되고, 이 점검은 강건성(조건을 바꿔도 결과가 버티는지)의 증거가 아니에요. 네 번째 점검은 엔도스랭크가 다시 출발 규칙에 얼마나 기대는지 물어요.

\dissertationSubsubheading[sub:damping]{감쇠 계수}

표~\ref{tab:robustness-damping}에는 같은 그래프에서 $d\in\{0.75,0.85,0.95\}$일 때 지표 가족 평균을 다시 계산한 값이 있어요.

% Generated by experiments/scripts/export_latex.py -- do not edit by hand.
% Source: experiments/data/2-processed-tables-and-evaluations/revised-usd/same-window/eval_summary.json
% Source: experiments/data/2-processed-tables-and-evaluations/revised-usd/benchmark/benchmark.json
\begin{table}[htbp]
\centering
\caption[강건성 점검: 페이지랭크 감쇠 계수 $d$에 따른 지표 가족 평균 켄달의 $\tau$.]{강건성 점검(조건을 바꿔도 결과가 버티는지 보기)이에요. 매칭 코호트(계약 $n=15{,}107$개)에서 페이지랭크 감쇠 계수 $d \in \{0.75, 0.85, 0.95\}$를 쓸 때 엔도스랭크(ER)와 AWP의 지표 가족 평균 켄달의 $\tau$를 보여 줘요. 열 이름에서 허락은 허락 가족, 송금은 송금 가족, 시빌은 시빌 안정성 가족이에요. 실행 시간은 전체 그래프에서 시간을 잰 풀이 프로그램 호출 다섯 번의 평균이에요. $d=0.85$의 실행 시간은 표~\ref{tab:benchmark-runtime}에서 잰 값을 다시 쓴 것이고, 새로 재지 않았어요(모든 실행은 표~\ref{tab:benchmark-runs}에 있어요). 신뢰구간(CI)은 계약을 다시 뽑아 구한 95\% 백분위 부트스트랩 구간이에요(짝지은 다시 뽑기 400번, 시드 42).}
\label{tab:robustness-damping}
\fitwidth{%
\begin{tabular}{rccccccrr}
\toprule
$d$ & ER 허락 & AWP 허락 & ER 송금 & AWP 송금 & ER 시빌 & AWP 시빌 & ER (초) & AWP (초) \\
\midrule
0.75 & \begin{tabular}[c]{@{}c@{}}0.586\\[-1pt]{\footnotesize [0.580, 0.594]}\end{tabular} & \begin{tabular}[c]{@{}c@{}}0.193\\[-1pt]{\footnotesize [0.182, 0.203]}\end{tabular} & \begin{tabular}[c]{@{}c@{}}0.158\\[-1pt]{\footnotesize [0.146, 0.168]}\end{tabular} & \begin{tabular}[c]{@{}c@{}}0.570\\[-1pt]{\footnotesize [0.562, 0.578]}\end{tabular} & \begin{tabular}[c]{@{}c@{}}0.007\\[-1pt]{\footnotesize [-0.002, 0.015]}\end{tabular} & \begin{tabular}[c]{@{}c@{}}0.397\\[-1pt]{\footnotesize [0.389, 0.406]}\end{tabular} & 5.984 & 15.236 \\
0.85 & \begin{tabular}[c]{@{}c@{}}0.585\\[-1pt]{\footnotesize [0.578, 0.592]}\end{tabular} & \begin{tabular}[c]{@{}c@{}}0.193\\[-1pt]{\footnotesize [0.182, 0.203]}\end{tabular} & \begin{tabular}[c]{@{}c@{}}0.158\\[-1pt]{\footnotesize [0.147, 0.169]}\end{tabular} & \begin{tabular}[c]{@{}c@{}}0.572\\[-1pt]{\footnotesize [0.564, 0.580]}\end{tabular} & \begin{tabular}[c]{@{}c@{}}0.007\\[-1pt]{\footnotesize [-0.001, 0.015]}\end{tabular} & \begin{tabular}[c]{@{}c@{}}0.395\\[-1pt]{\footnotesize [0.387, 0.403]}\end{tabular} & 10.294 & 22.367 \\
0.95 & \begin{tabular}[c]{@{}c@{}}0.583\\[-1pt]{\footnotesize [0.577, 0.591]}\end{tabular} & \begin{tabular}[c]{@{}c@{}}0.192\\[-1pt]{\footnotesize [0.181, 0.202]}\end{tabular} & \begin{tabular}[c]{@{}c@{}}0.158\\[-1pt]{\footnotesize [0.147, 0.169]}\end{tabular} & \begin{tabular}[c]{@{}c@{}}0.577\\[-1pt]{\footnotesize [0.569, 0.585]}\end{tabular} & \begin{tabular}[c]{@{}c@{}}0.008\\[-1pt]{\footnotesize [-0.001, 0.016]}\end{tabular} & \begin{tabular}[c]{@{}c@{}}0.391\\[-1pt]{\footnotesize [0.384, 0.400]}\end{tabular} & 31.233 & 60.962 \\
\bottomrule
\end{tabular}
}
\end{table}


엔도스랭크에서는 세 값 사이에서 어떤 지표 가족 평균도 $0.003$보다 더 움직이지 않았어요. 엔도스랭크의 허락 계수는 $d=0.75$에서 $0.586$ $[0.580,0.594]$이고, $d=0.95$에서 $0.583$ $[0.577,0.591]$이에요. AWP는 송금 가족에서 가장 많이 움직였어요. $0.570$ $[0.562,0.578]$에서 $0.577$ $[0.569,0.585]$로 바뀌었어요. 그리고 AWP의 시빌 안정성 계수는 $0.397$ $[0.389,0.406]$에서 $0.391$ $[0.384,0.400]$로 내려갔어요. 두 점수 어느 쪽의 모양도, 흔히 쓰는 $d=0.85$ 때문에 생긴 겉모습이 아니에요. 엔도스랭크가 잘 버티는 것은 대부분 그 구조 때문이에요. 주인들이 아무 허락도 받지 않은 스펜더에서는, 엔도스랭크가 $c\,(1+d\,\delta(v))$와 같아요. 여기서 $\delta(v)$는 주인 기준으로 나누어 맞춘(정규화한) 들어오는 무게의 합이에요(\ref{sub:theory-closed-form}절). 이런 스펜더들 사이에서 이 식이 주는 순서는 $d$에 달려 있지 않아요. 들어오는 화살표가 있는 코호트 계약 14,406개 가운데 12,039개가 이런 스펜더예요. 반복 횟수는 $d$가 커질수록 늘어났어요. 엔도스랭크는 41번에서 69번, 217번으로, AWP는 55번에서 92번, 261번으로 늘었어요. 실행 시간은 엔도스랭크가 5.98, 10.29, 31.23~초이고, AWP가 15.24, 22.37, 60.96~초예요. 시간 비율은 2.55, 2.17, 1.95이고, 이때 화살표 수 비율은 2.23이에요. 이 시간은 표~\ref{tab:benchmark-runtime}의 전체 그래프에서 잰 거예요(엔도스랭크는 점 6,219,063개와 화살표 11,210,114개, AWP는 7,882,356개와 25,001,216개). 그래프와 감쇠 계수의 짝마다 시간을 한 번씩 쟀어요. $d=0.85$의 값은 주된 측정값을 다시 쓴 거예요. 그림~\ref{fig:damping-sensitivity}에는 지표 가족 평균과 그 구간이 그려져 있어요.

\begin{figure}[htbp]
\centering
\includegraphics[width=0.85\textwidth]{\figpath{damping-sensitivity.pdf}}
\caption[{$d\in\{0.75,0.85,0.95\}$에서 지표 가족 평균 켄달의 $\tau$의 민감도(표~\ref{tab:robustness-damping}의 값).}]{엔도스랭크와 AWP의 지표 가족 평균 켄달의 $\tau$가 감쇠 계수에 따라 얼마나 바뀌는지 보여 줘요. 송금, 허락, 시빌 안정성 가족에서 $d\in\{0.75,0.85,0.95\}$에 걸쳐 그렸고, 95\% 부트스트랩 구간도 함께 그렸어요(값은 표~\ref{tab:robustness-damping}에서 가져왔어요. 매칭 코호트, 계약 $n=15{,}107$개).}
\label{fig:damping-sensitivity}
\end{figure}

\dissertationSubsubheading[sub:token-subgraph]{금액 기울기 민감도}

표~\ref{tab:robustness-slope}에는 금액 무게의 기울기 $b$를 10으로 나누거나 10을 곱해서(\ref{sub:top-tokens}절), $T_{\mathrm{obs}}$에서 엔도스랭크와 AWP를 다시 계산한 값이 있어요. 그래프와 대리 지표는 똑같이 두었어요.

% Generated by experiments/scripts/export_latex.py -- do not edit by hand.
% Source: experiments/data/2-processed-tables-and-evaluations/revised-usd/robustness/robustness.json
% Source: experiments/data/0-usd-token-list-and-prices/constants.json
\begin{table}[htbp]
\centering
\caption[강건성 점검: 금액 기울기를 열 배 작게 하거나 크게 했을 때의 지표 가족 평균 켄달의 $\tau$.]{강건성 점검이에요. 금액 무게 $V$의 기울기 $b$를 열로 나누거나 열 배로 키웠을 때, 매칭 코호트(계약 $n=15{,}107$개)의 지표 가족 평균 켄달의 $\tau$를 주 분석 옆에 보여 줘요. 주 분석의 기울기는 금액이 중앙값일 때 $V=1/2$이 되게 정했어요. 그 중앙값은 송금이 \$200, 허락이 \$62예요. $b/10$은 그 점을 중앙값의 열 배로 옮기고, $10b$는 중앙값의 십분의 일로 옮겨요. 각 변형과, 같은 방법(엔도스랭크와 AWP)의 주 점수 사이의 순위 일치 $\tau$는 $b/10$에서 $0.973$과 $0.973$이고, $10b$에서 $0.968$과 $0.971$이에요. 변형에 대해서는 구간을 계산하지 않았어요.}
\label{tab:robustness-slope}
\fitwidth{%
\begin{tabular}{lccccccccc}
\toprule
방법 & \multicolumn{3}{c}{송금} & \multicolumn{3}{c}{허락} & \multicolumn{3}{c}{시빌} \\
\cmidrule(lr){2-4} \cmidrule(lr){5-7} \cmidrule(lr){8-10}
 & $b/10$ & $b$ & $10b$ & $b/10$ & $b$ & $10b$ & $b/10$ & $b$ & $10b$ \\
\midrule
EndorseRank & 0.157 & 0.158 & 0.159 & 0.588 & 0.585 & 0.580 & 0.006 & 0.007 & 0.009 \\
AWP & 0.577 & 0.572 & 0.567 & 0.193 & 0.193 & 0.191 & 0.397 & 0.395 & 0.393 \\
\bottomrule
\end{tabular}
}
\end{table}


어떤 지표 가족 평균도 $0.0054$보다 더 움직이지 않았고, 모양도 그대로였어요. 허락 가족에서는 엔도스랭크가 앞서고, 송금 가족과 시빌 안정성 가족에서는 AWP가 앞서요. 더 가파른 기울기 $10b$에서는 무게가 2분의 1이 되는 금액이 중앙값 금액의 10분의 1이 돼요. 그래서 대부분의 금액이 1에 가까운 무게를 받고, 화살표 무게는 이벤트를 센 수에 더 가까워져요. 이 기울기는 엔도스랭크의 허락 계수를 $0.585$에서 $0.580$으로, AWP의 송금 계수를 $0.572$에서 $0.567$로 낮춰요. 더 완만한 기울기 $b/10$은 이 계수들을 $0.588$과 $0.577$로 높여요. 각 변형은 같은 방법의 주 점수와 거의 같은 순서로 계약을 세워요. $b/10$에서는 두 점수 모두 $\tau=0.973$이고, $10b$에서는 엔도스랭크가 $0.968$, AWP가 $0.971$이에요. 변형들에 대해서는 구간을 계산하지 않았어요. 그래서 가족 수준의 결과는, 금액 무게가 2분의 1이 되는 금액이 얼마인지에 달려 있지 않아요.

\dissertationSubsubheading[sub:sample-size]{표본 정의 민감도}

표~\ref{tab:robustness-sample-size}에는 \ref{sub:cohorts}절의 부분 표본에서 지표 가족 평균을 다시 계산한 값이 있어요. 부분 표본마다 그 표본의 계약에 닿는 화살표만으로 점수를 매겼어요.

% Generated by experiments/scripts/export_latex.py -- do not edit by hand.
% Source: experiments/data/2-processed-tables-and-evaluations/revised-usd/robustness/robustness.json
% Source: experiments/data/2-processed-tables-and-evaluations/revised-usd/benchmark/benchmark.json
\begin{table}[htbp]
\centering
\caption[평가 집합에 따른 지표 가족 평균의 민감도.]{평가 집합이 바뀌면 지표 가족 평균이 얼마나 움직이는지 보여 줘요. 매칭 코호트에서 SHA256으로 늘 같은 결과가 나오게 뽑은 부분 표본마다(시드 \texttt{contract-benchmark-v1}, 코호트 전체는 계약 $n=15{,}107$개) 엔도스랭크(ER)와 AWP의 지표 가족 평균 켄달의 $\tau$를 구했어요. 단계마다 그 단계의 계약에 닿는 화살표로 점수를 매겼어요. 열 이름에서 “허락”은 허락, “송금”은 송금, “시빌”은 시빌 안정성을 뜻해요. 행마다 평가한 계약 집합이 달라요. 그래서 이 $\tau$ 값들은 표~\ref{tab:alignment-hybrid}의 값과 비교할 수 없어요. 이 값들은 평가하는 집단에 따라 통계량이 얼마나 움직이는지 보여 줘요. 단계마다 구간은 계산하지 않았어요. 실행 시간은 표~\ref{tab:benchmark-scaling}의 측정값이에요.}
\label{tab:robustness-sample-size}
\fitwidth{%
\begin{tabular}{rccccccrrr}
\toprule
$n$ & ER 허락 & AWP 허락 & ER 송금 & AWP 송금 & ER 시빌 & AWP 시빌 & ER (초) & AWP (초) & AWP/ER 시간 \\
\midrule
1{,}000 & 0.715 & 0.196 & 0.205 & 0.586 & 0.038 & 0.408 & 0.580 & 1.677 & 2.9$\times$ \\
2{,}000 & 0.686 & 0.199 & 0.196 & 0.588 & 0.033 & 0.411 & 0.939 & 4.272 & 4.6$\times$ \\
5{,}000 & 0.634 & 0.211 & 0.195 & 0.571 & 0.034 & 0.391 & 2.496 & 8.995 & 3.6$\times$ \\
10{,}000 & 0.607 & 0.196 & 0.167 & 0.571 & 0.013 & 0.394 & 6.220 & 16.821 & 2.7$\times$ \\
15{,}107 & 0.585 & 0.193 & 0.158 & 0.572 & 0.007 & 0.395 & 10.294 & 22.367 & 2.2$\times$ \\
\bottomrule
\end{tabular}
}
\end{table}


계수는 표본에 따라 움직여요. 엔도스랭크가 가장 많이 움직여요. 엔도스랭크의 허락 계수는 $n=1{,}000$에서 $0.715$이고, 단계마다 내려가서 전체 코호트에서는 $0.585$가 돼요. 엔도스랭크의 송금 계수는 $0.205$에서 $0.158$로 내려가요. AWP는 모든 가족에서 전체 코호트 값과의 차이가 $0.018$ 안에 머물러요. 단계들에 대해서는 구간을 계산하지 않았어요. 이런 변화 가운데 어느 것도 점수에 대한 증거가 아니에요. 한 단계는 뽑힌 계약에 닿는 화살표만 남겨요. 그래서 뽑히지 않은 코호트 계약을 지나는 길은 끊겨요. 그러면 각 점수는 뽑힌 계약 자신의 화살표를 센 수에 더 가까워지고, 그것이 바로 대리 지표가 세는 것이에요. 그래서 이 연구는 이 표를 표본 정의에 대한 민감도 분석으로 보고해요. 이 표의 값을 표~\ref{tab:alignment-hybrid}의 값과 견주어서는 안 돼요. 주된 표본은 매칭 코호트예요.

\dissertationSubsubheading[sub:endorserank-restarts]{엔도스랭크의 다시 출발}

엔도스랭크는 걷기를 고르게 다시 출발시켜요. 이와 달리 AWP는 각 주소가 최근에 보낸 활동에 비례해서 그 주소에서 걷기를 다시 출발시켜요(\ref{sub:endorserank-graph}절). AWP의 규칙을 허락에 옮겨 쓰면, 주인이 허락해 준 스펜더 수에 비례해서 그 주인에서 걷기를 다시 출발시켜요. 이때 스펜더마다 마지막 허락이 얼마나 오래되었는지로 무게를 매겨요. 분석 계획은 엔도스랭크에 고르게 다시 출발하기를 정했어요. 그리고 AWP의 다시 출발 규칙을 쓴 변형을 그 옆에 함께 보고할 점수로 정했어요. 하지만 둘 사이의 비교는 하나도 정하지 않았어요. 표~\ref{tab:endorserank-restarts}에는 두 버전의 점수가 있고, 그 옆에 AWP도 있어요.

% Generated by experiments/scripts/export_latex.py -- do not edit by hand.
% Source: experiments/data/2-processed-tables-and-evaluations/revised-usd/same-window/eval_summary.json
% Source: experiments/data/2-processed-tables-and-evaluations/revised-usd/holdout-w0/eval_summary.json
% Source: experiments/data/2-processed-tables-and-evaluations/revised-usd/window-w1/eval_summary.json
\begin{table}[htbp]
\centering
\caption[AWP 옆에 놓은, 고르게 다시 출발하는 엔도스랭크와 활동에 따라 무게를 단 재시작을 쓰는 엔도스랭크.]{엔도스랭크를 고르게 다시 출발하게 한 것과, 허락하는 활동에 따라 무게를 단 재시작으로 다시 출발하게 한 것을 AWP 옆에 놓았어요. 뒤의 재시작은 AWP의 규칙을 허락 그래프에 쓴 것이에요. 매칭 코호트(계약 $n=15{,}107$개)의 지표 가족 평균 켄달의 $\tau$를 보여 줘요. 또 두 스펜더 라벨과의 $\tau$를 W0에서(스펜더 계약 $n=17{,}050$개, 점수는 2026년 3월 31일에 동결), 그리고 W1에서($n=18{,}152$, 2026년 6월 30일에 동결) 보여 줘요. 두 엔도스랭크 행 모두 각 마지막 허락에 $\sigma(\Delta t)V(z)$만큼 무게를 줘요. 신뢰구간(CI)은 계약을 다시 뽑아 구한 95\% 백분위 부트스트랩 구간이에요(짝지은 다시 뽑기 400번, 시드 42).}
\label{tab:endorserank-restarts}
\fitwidth{%
\begin{tabular}{llccccccc}
\toprule
 & & \multicolumn{3}{c}{같은 기간, 매칭 코호트} & \multicolumn{2}{c}{W0, 스펜더} & \multicolumn{2}{c}{W1, 스펜더} \\
\cmidrule(lr){3-5} \cmidrule(lr){6-7} \cmidrule(lr){8-9}
방법 & 재시작 & 송금 & 허락 & 시빌 & 새 허락 쌍 & 새 송금자 & 새 허락 쌍 & 새 송금자 \\
\midrule
EndorseRank & 고르게 & \begin{tabular}[c]{@{}c@{}}0.158\\[-1pt]{\footnotesize [0.147, 0.169]}\end{tabular} & \begin{tabular}[c]{@{}c@{}}0.585\\[-1pt]{\footnotesize [0.578, 0.592]}\end{tabular} & \begin{tabular}[c]{@{}c@{}}0.007\\[-1pt]{\footnotesize [-0.001, 0.015]}\end{tabular} & \begin{tabular}[c]{@{}c@{}}0.231\\[-1pt]{\footnotesize [0.220, 0.242]}\end{tabular} & \begin{tabular}[c]{@{}c@{}}0.200\\[-1pt]{\footnotesize [0.189, 0.210]}\end{tabular} & \begin{tabular}[c]{@{}c@{}}0.234\\[-1pt]{\footnotesize [0.222, 0.244]}\end{tabular} & \begin{tabular}[c]{@{}c@{}}0.194\\[-1pt]{\footnotesize [0.184, 0.205]}\end{tabular} \\
EndorseRank & 허락하는 활동에 따라 & \begin{tabular}[c]{@{}c@{}}0.164\\[-1pt]{\footnotesize [0.153, 0.173]}\end{tabular} & \begin{tabular}[c]{@{}c@{}}0.426\\[-1pt]{\footnotesize [0.418, 0.435]}\end{tabular} & \begin{tabular}[c]{@{}c@{}}0.049\\[-1pt]{\footnotesize [0.039, 0.057]}\end{tabular} & \begin{tabular}[c]{@{}c@{}}0.308\\[-1pt]{\footnotesize [0.299, 0.319]}\end{tabular} & \begin{tabular}[c]{@{}c@{}}0.269\\[-1pt]{\footnotesize [0.260, 0.278]}\end{tabular} & \begin{tabular}[c]{@{}c@{}}0.295\\[-1pt]{\footnotesize [0.285, 0.304]}\end{tabular} & \begin{tabular}[c]{@{}c@{}}0.242\\[-1pt]{\footnotesize [0.232, 0.253]}\end{tabular} \\
AWP & 송금하는 활동에 따라 & \begin{tabular}[c]{@{}c@{}}0.572\\[-1pt]{\footnotesize [0.564, 0.580]}\end{tabular} & \begin{tabular}[c]{@{}c@{}}0.193\\[-1pt]{\footnotesize [0.182, 0.203]}\end{tabular} & \begin{tabular}[c]{@{}c@{}}0.395\\[-1pt]{\footnotesize [0.387, 0.403]}\end{tabular} & \begin{tabular}[c]{@{}c@{}}0.196\\[-1pt]{\footnotesize [0.181, 0.210]}\end{tabular} & \begin{tabular}[c]{@{}c@{}}0.346\\[-1pt]{\footnotesize [0.336, 0.356]}\end{tabular} & \begin{tabular}[c]{@{}c@{}}0.169\\[-1pt]{\footnotesize [0.157, 0.183]}\end{tabular} & \begin{tabular}[c]{@{}c@{}}0.330\\[-1pt]{\footnotesize [0.321, 0.340]}\end{tabular} \\
\bottomrule
\end{tabular}
}
\end{table}


활동 재시작(활발한 주소에서 더 자주 다시 출발하기)을 쓰면, 엔도스랭크의 같은 기간 허락 계수는 $0.585$ $[0.578,0.592]$에서 $0.426$ $[0.418,0.435]$로 내려가요. 송금 계수는 거의 바뀌지 않아요($0.164$ $[0.153,0.173]$ 대 $0.158$ $[0.147,0.169]$). 이때 다시 출발하는 몫은 허락을 가장 많이 하는 주인들에게 가요. 그래서 이 점수는 스펜더의 주인이 몇 명인지뿐만 아니라, 그 주인들이 얼마나 활발한지로도 스펜더에 무게를 매겨요. 기간 밖에서는 이 변형이 두 기간 모두에서 두 스펜더 라벨을 더 잘 줄 세워요. W0에서 이 변형은 새 허락 쌍(새로 생긴 주인--스펜더 허락)에서 $0.308$ $[0.299,0.319]$, 새 송금자(새로 송금을 보내온 주소)에서 $0.269$ $[0.260,0.278]$에 이르러요. 이것은 고르게 다시 출발하는 버전보다 $+0.077$ $[0.071,0.084]$와 $+0.070$ $[0.064,0.076]$만큼 높아요. W1에서는 $0.295$ $[0.285,0.304]$와 $0.242$ $[0.232,0.253]$에 이르러요. 이것은 고르게 다시 출발하는 버전보다 $+0.061$ $[0.054,0.067]$과 $+0.048$ $[0.042,0.053]$만큼 높아요. 이 대비들은 각 기간의 라벨을 본 뒤에 계산했고, 어떤 규칙에도 들어 있지 않아요. 고르게 다시 출발하는 규칙은 같은 기간에서 허락 구성개념을 지켜요. 활동 규칙은 나중에 들어오는 것을 더 잘 줄 세워요. W0에서 이 변형은 새 허락 쌍을 받은 허락 수보다 더 잘 줄 세워요($+0.059$, $[0.050,0.067]$). 반면에 고르게 다시 출발하는 버전은 받은 허락 수보다 못해요(\ref{sec:holdout-results}절).

\dissertationSubheading[sec:holdout-results]{시간 홀드아웃(RQ4, RQ5)}

같은 기간 점검으로는 점수와 그 점수 자신의 구성개념을 떼어 놓을 수 없어요. 그 대신 홀드아웃은, 어떤 기준 시점까지 잰 중심성(그래프에서 얼마나 중요한 자리에 있는지)이 그 뒤의 행동을 예측하는지(미리 맞히는지) 시험해요. \ref{sec:holdout-protocol}절에 따라, W0 스펜더 코호트의 점수를 $t_1$에 동결해요. 그다음 2026년 4월 1일부터 6월 30일까지 계약마다 두 라벨을 세요. 하나는 새 주인--스펜더 허락 쌍(보증 구성개념)이고, 다른 하나는 새 송금자(흐름 구성개념)예요. 표~\ref{tab:holdout-spenders}에서는 $t_1$에 계산한 모든 점수를 두 라벨과 맞대어 봐요. 두 원시 차수 기준선도 함께 넣었어요. 표~\ref{tab:holdout-diff}에는 짝지은 대비가 있어요. 첫 번째 묶음에는 규칙~A의 두 대비가 있어요. 규칙~A는 이 연구의 점수가 하나도 없을 때 분석 계획에서 정했고, 여기서는 고친 점수에 적용했어요. 두 번째 묶음에는 라벨을 본 뒤에 계산한 대비들이 있어요.

% Generated by experiments/scripts/export_latex.py -- do not edit by hand.
% Source: experiments/data/2-processed-tables-and-evaluations/revised-usd/holdout-w0/eval_summary.json
\begin{table}[htbp]
\centering
\caption[스펜더 코호트의 시간 홀드아웃 W0.]{스펜더 코호트(계약 $n=17{,}050$개)에서 해 본 시간 홀드아웃 W0이에요. 2026년 3월 31일에 동결한 점수와, 2026년 4월 1일부터 2026년 6월 30일까지 그 계약에서 센 두 라벨 사이의 켄달의 $\tau$를 보여 줘요. 두 라벨은 새로 생긴 주인--스펜더 허락 쌍의 수와, 그 계약에 처음으로 송금을 보낸 주소의 수예요. 코호트에는 동결일에 마지막 허락이 0보다 큰 스펜더 17{,}167개 가운데 계약 17{,}050개를 남겼어요. 코드가 없는 EOA 104개와 EIP-7702 위임이 있는 EOA 13개는 뺐어요. 계약 2{,}070개는 새 쌍을 하나 이상 받았고, 1{,}695개는 새 송금자를 하나 이상 받았어요. 기준선 행은 각 한 층 페이지랭크가 매끄럽게 다듬는 $t_1$의 원시 차수예요. C-PR ($\lambda=1$)과 C-PR ($\lambda=0$)은 결합 연산자의 허락 층과 송금 층을 각각 한 층만 걷는 것이에요. 신뢰구간(CI)은 계약을 다시 뽑아 구한 95\% 백분위 부트스트랩 구간이에요(짝지은 다시 뽑기 400번, 시드 42).}
\label{tab:holdout-spenders}
\fitwidth{%
\begin{tabular}{lcc}
\toprule
$t_1$의 점수 & 새 허락 쌍 ($\tau$ [95\% 신뢰구간]) & 새 송금자 ($\tau$ [95\% 신뢰구간]) \\
\midrule
$t_1$의 받은 허락 수 (기준선) & 0.250 [0.237, 0.262] & 0.180 [0.169, 0.192] \\
$t_1$의 송금 입차수 (기준선) & 0.135 [0.121, 0.150] & 0.303 [0.292, 0.315] \\
\midrule
EndorseRank & 0.231 [0.220, 0.242] & 0.200 [0.189, 0.210] \\
AWP & 0.196 [0.181, 0.210] & 0.346 [0.336, 0.356] \\
\midrule
C-PR ($\lambda=1$) & 0.223 [0.213, 0.233] & 0.188 [0.177, 0.199] \\
C-PR ($\lambda=0$) & 0.160 [0.144, 0.174] & 0.296 [0.285, 0.307] \\
C-PR ($\lambda=0.25$) & 0.253 [0.243, 0.264] & 0.269 [0.257, 0.280] \\
C-PR ($\lambda=0.5$) & 0.252 [0.241, 0.262] & 0.262 [0.250, 0.273] \\
C-PR ($\lambda=0.75$) & 0.249 [0.238, 0.260] & 0.253 [0.242, 0.265] \\
S-PR & 0.172 [0.157, 0.185] & 0.295 [0.284, 0.305] \\
\bottomrule
\end{tabular}
}
\end{table}


% Generated by experiments/scripts/export_latex.py -- do not edit by hand.
% Source: experiments/data/2-processed-tables-and-evaluations/revised-usd/holdout-w0/eval_summary.json
\begin{table}[htbp]
\centering
\caption[스펜더 홀드아웃 W0에서 짝지은 부트스트랩으로 구한 차이.]{스펜더 홀드아웃 W0에서 짝지은 부트스트랩으로 구한 차이예요(계약 $n=17{,}050$개, $a$와 $b$가 함께 쓰는 다시 뽑은 표본 400개, 시드 42). 첫 번째 묶음에는 규칙 A의 주된 부분에 들어가는 두 대비가 있어요. 이 두 대비는 어떤 자료도 보기 전인 2026년 10월 9일에 커밋한 분석 계획에서 미리 정했고, 여기서는 고친 점수에 적용했어요. 하나는 새 허락 쌍에서 C-PR과 C-PR ($\lambda=1$)을 비교하고, 다른 하나는 새 송금자에서 C-PR과 C-PR ($\lambda=0$)을 비교해요. 두 번째 묶음은 라벨을 본 뒤에 계산했어요. 여기에는 다음이 들어 있어요. 허락 층만 걷는 C-PR ($\lambda=1$)이 $t_1$에서의 자기 원시 차수보다 더 얻은 값(증가분), S-PR에 대한 같은 비교들, 엔도스랭크와 AWP가 각자의 원시 차수보다 더 얻은 값, 새 허락 쌍에서 C-PR과 C-PR ($\lambda=0$)의 비교, 그리고 C-PR과 엔도스랭크, C-PR과 AWP의 비교예요. 규칙 A의 주된 부분은 C-PR이 새 허락 쌍에서 C-PR ($\lambda=1$)보다 더 나쁘지 않고, 새 송금자에서 C-PR ($\lambda=0$)보다 더 나아야 한다는 것이에요. “더 나쁘지 않음”은 성립하지만 “더 나음”은 성립하지 않아요. 그래서 이 규칙은 충족되지 않아요. 엔도스랭크와 AWP 사이의 차이는 표~\ref{tab:endorserank-awp-holdout}에 있어요.}
\label{tab:holdout-diff}
\fitwidth{%
\begin{tabular}{llccccc}
\toprule
대비 ($a$ 빼기 $b$) & 라벨 & $\tau_a$ & $\tau_b$ & $\Delta\tau$ & $\Delta\tau$의 95\% 신뢰구간 & $P(\Delta\tau>0)$ \\
\midrule
\multicolumn{7}{l}{\emph{규칙 A, 주된 부분, 계획에서 미리 정함}} \\
C-PR 빼기 C-PR ($\lambda=1$) & 새 허락 쌍 & 0.252 & 0.223 & +0.029 & [0.022, 0.036] & 1.000 \\
C-PR 빼기 C-PR ($\lambda=0$) & 새 송금자 & 0.262 & 0.296 & -0.034 & [-0.040, -0.029] & 0.000 \\
\midrule
\multicolumn{7}{l}{\emph{라벨을 본 뒤에 계산함}} \\
C-PR ($\lambda=1$) 빼기 $t_1$ 받은 허락 수 & 새 허락 쌍 & 0.223 & 0.250 & -0.027 & [-0.034, -0.019] & 0.000 \\
S-PR 빼기 C-PR ($\lambda=1$) & 새 허락 쌍 & 0.172 & 0.223 & -0.051 & [-0.066, -0.038] & 0.000 \\
S-PR 빼기 C-PR ($\lambda=0$) & 새 송금자 & 0.295 & 0.296 & -0.001 & [-0.004, 0.002] & 0.265 \\
EndorseRank 빼기 $t_1$ 받은 허락 수 & 새 허락 쌍 & 0.231 & 0.250 & -0.019 & [-0.026, -0.011] & 0.000 \\
AWP 빼기 $t_1$ 입차수 & 새 송금자 & 0.346 & 0.303 & +0.043 & [0.038, 0.048] & 1.000 \\
C-PR 빼기 C-PR ($\lambda=0$) & 새 허락 쌍 & 0.252 & 0.160 & +0.092 & [0.080, 0.105] & 1.000 \\
C-PR 빼기 EndorseRank & 새 허락 쌍 & 0.252 & 0.231 & +0.021 & [0.013, 0.028] & 1.000 \\
C-PR 빼기 AWP & 새 송금자 & 0.262 & 0.346 & -0.085 & [-0.092, -0.078] & 0.000 \\
C-PR 빼기 AWP & 새 허락 쌍 & 0.252 & 0.196 & +0.056 & [0.042, 0.070] & 1.000 \\
\bottomrule
\end{tabular}
}
\end{table}


첫 번째 발견은 RQ4의 두 번째 부분에 답해요. 그리고 그 답은 두 걷기에서 서로 달라요. 송금 쪽에서 AWP는 새 송금자를 $\tau=0.346$ $[0.336,0.356]$으로 줄 세워요. 이것은 $t_1$의 송금 입차수인 $0.303$ $[0.292,0.315]$보다 높고, 짝지은 차이는 $+0.043$ $[0.038,0.048]$이에요. 그러니까 AWP의 무게와 활동 재시작으로 퍼뜨리는 일(전파)이, 이 라벨에서 그냥 센 수에 정보를 보탰어요. 달러 금액과 고르게 다시 출발하기로 걷는 송금 층은 이렇게 뚜렷하게 보태지 못해요. C-PR ($\lambda=0$)은 $0.296$ $[0.285,0.307]$에 이르러, 입차수보다 조금 낮아요(짝지은 대비는 계산하지 않았어요). 허락 쪽에서는 그냥 센 수가 이겨요. $t_1$의 받은 허락 수는 새 허락 쌍을 $0.250$ $[0.237,0.262]$로 줄 세워요. 엔도스랭크는 $0.231$ $[0.220,0.242]$이고(짝지은 차이 $-0.019$, $[-0.026,-0.011]$), C-PR ($\lambda=1$)은 $0.223$ $[0.213,0.233]$이에요($-0.027$, $[-0.034,-0.019]$). 스펜더는 이미 가진 주인이 많을수록 새 주인을 더 많이 얻어요. 선호적 연결(이미 많이 가진 쪽이 더 많이 얻는 현상)이 예측하는 그대로예요\parencite{barabasi1999scaling}. 그냥 센 수는 이것을 바로 잡아내요. 엔도스랭크는 각 주인의 보증을, 그 주인이 허락한 모든 스펜더에게 나눠요(\ref{sub:theory-closed-form}절). 그래서 주인들이 다른 스펜더도 많이 허락한 스펜더는, 그냥 센 수만으로 놓일 자리보다 낮게 줄 세워져요. 이 라벨에서는 그렇게 나누는 것이 조금 손해예요. 이런 차수 대비 가운데 어느 것도 계획에서 정하지 않았어요. 모두 라벨을 본 뒤에 계산했어요.

표 전체를 가로질러 읽으면, 각 라벨은 자기와 같은 종류의 화살표를 쓰는 점수가 가장 잘 줄 세워요. 새 허락 쌍에서는 C-PR 혼합 점수들, 받은 허락 수, 엔도스랭크가 앞서요. AWP는 $0.196$ $[0.181,0.210]$, S-PR은 $0.172$ $[0.157,0.185]$, C-PR ($\lambda=0$)은 $0.160$ $[0.144,0.174]$, 송금 입차수는 $0.135$ $[0.121,0.150]$로 더 낮아요. 새 송금자에서는 AWP와 송금 입차수가 앞서요. 허락 쪽 점수들은 더 낮고, 서로 가까워요. 엔도스랭크는 $0.200$ $[0.189,0.210]$, C-PR ($\lambda=1$)은 $0.188$ $[0.177,0.199]$, 받은 허락 수는 $0.180$ $[0.169,0.192]$에 이르러요. 이들의 구간은 서로 겹치고, 이들 사이의 짝지은 대비는 없어요. 여기서 엔도스랭크는 동점을 거의 남기지 않아요. 엔도스랭크는 스펜더 17,050개에 서로 다른 점수 16,208개를 줘요. 가장 큰 동점 덩어리에는 그중 156개가 들어 있고, 쌍의 0.027\%만 동점이에요. 그래서 엔도스랭크의 계수는 거의 모든 스펜더의 순서에 바탕을 두어요.

어떤 계약이 그래프에 들어 있는지(포함 범위)로는 차수 비교를 설명할 수 없어요. 스펜더 17,050개 가운데 2,891개는 동결일까지 송금 화살표가 하나도 없어서, AWP에서 0점을 받아요. 이 스펜더들은 나머지보다 새 허락 쌍을 조금 더 자주 얻었고(13.3\% 대 11.9\%), 새 송금자는 거의 얻지 못했어요(0.6\% 대 11.9\%). 표~\ref{tab:holdout-receiving}에서는 홀드아웃을, 동결 전에 다른 주소에게서 송금을 받은 적이 있는 스펜더 14,094개로 좁혀요. 이렇게 좁히기로 한 것은 결과를 안 뒤에 정했어요. 여기서도 AWP는 여전히 새 송금자를 송금 입차수보다 더 잘 줄 세워요($+0.056$, $[0.049,0.064]$). 그리고 엔도스랭크는 여전히 새 허락 쌍을 받은 허락 수보다 못하게 줄 세워요($-0.013$, $[-0.020,-0.006]$). 이 부분 집합에서 새 허락 쌍에 대한 AWP의 계수는 $0.196$ $[0.181,0.210]$에서 $0.298$ $[0.287,0.308]$로 올라가요. AWP의 0점들이, 나중에 허락을 얻은 많은 스펜더를 맨 아래에 놓았었기 때문이에요. 이 부분 집합에서 AWP의 점 추정값은 받은 허락 수($0.249$ $[0.237,0.264]$)와 엔도스랭크($0.236$ $[0.224,0.249]$)의 점 추정값보다 높아요. 다만 이 쌍들에 대해서는 짝지은 대비를 계산하지 않았어요.

% Generated by experiments/scripts/export_latex.py -- do not edit by hand.
% Source: experiments/data/2-processed-tables-and-evaluations/revised-usd/holdout-w0/eval_summary.json
\begin{table}[htbp]
\centering
\caption[동결일 전에 송금을 받은 적이 있는 스펜더로 좁힌 스펜더 홀드아웃 W0.]{스펜더 홀드아웃 W0을 좁혀서 본 결과예요. 스펜더 계약 17{,}050개 가운데, 동결일 전에 다른 주소에서 송금을 받은 적이 있는 14{,}094개만 남겼어요($n=14{,}094$). 나머지 2{,}956개에는 다른 주소에서 들어오는 송금 화살표가 없어요. 그리고 스펜더 2{,}891개는 AWP 점수가 0이에요. 이렇게 좁히는 방법은 결과를 안 뒤에 골랐어요. 신뢰구간(CI)은 계약을 다시 뽑아 구한 95\% 백분위 부트스트랩 구간이에요(짝지은 다시 뽑기 400번, 시드 42).}
\label{tab:holdout-receiving}
\fitwidth{%
\begin{tabular}{lcc}
\toprule
$t_1$의 점수 & 새 허락 쌍 ($\tau$ [95\% 신뢰구간]) & 새 송금자 ($\tau$ [95\% 신뢰구간]) \\
\midrule
EndorseRank & 0.236 [0.224, 0.249] & 0.231 [0.219, 0.244] \\
AWP & 0.298 [0.287, 0.308] & 0.357 [0.347, 0.367] \\
\midrule
$t_1$의 받은 허락 수 & 0.249 [0.237, 0.264] & 0.205 [0.193, 0.218] \\
$t_1$의 송금 입차수 & 0.211 [0.196, 0.224] & 0.301 [0.289, 0.313] \\
\midrule
EndorseRank 빼기 $t_1$ 받은 허락 수 (새 허락 쌍) & \multicolumn{2}{c}{$\Delta\tau=-0.013$ [-0.020, -0.006]} \\
AWP 빼기 $t_1$ 입차수 (새 송금자) & \multicolumn{2}{c}{$\Delta\tau=+0.056$ [0.049, 0.064]} \\
\bottomrule
\end{tabular}
}
\end{table}


두 번째 발견은, 고친 점수에서 결합 연산자가 규칙~A를 통과하지 못한다는 거예요. 다만 탐색적(미리 정하지 않고 살펴본) 비교에서는, 결합 연산자가 자기 송금 층만 걷는 것보다 나아요. 규칙~A는 이 연구의 점수가 하나도 없을 때 분석 계획에서 정했어요. 규칙~A는 $\lambda=0.5$인 C-PR에게 두 가지를 요구했어요. 새 허락 쌍에서는 C-PR ($\lambda=1$)보다 나쁘지 않아야 하고, 새 송금자에서는 C-PR ($\lambda=0$)보다 나아야 해요. 여기서는 등록 분석의 결과를 안 뒤에 정의한 점수에 이 규칙을 적용해요. 그래서 그 판정은 재분석이고, 등록한 시험이 아니에요. 첫 번째 조건은 성립해요. C-PR은 허락 걷기보다 위에 있어요($+0.029$, $[0.022,0.036]$). 두 번째 조건은 성립하지 않아요. C-PR은 새 송금자에서 송금 걷기보다 아래에 있어요($-0.034$, $[-0.040,-0.029]$). 등록 분석도 같은 판정에 이르렀어요(표~\ref{tab:registered-vs-revised}). 주된 부분이 실패했으므로, RQ5의 뒷부분은 계획에서 정한 대로 “아니다”라는 답을 받아요. RQ5의 뒷부분은 C-PR이 어느 한 층만 걷는 걷기를 이기는지 묻는 거예요. 규칙~A에 들어 있지 않은 비교로, 새 허락 쌍에서 송금 걷기 하나와 견주면 C-PR이 더 나아요. C-PR은 $0.252$ $[0.241,0.262]$에 이르고, 송금 걷기는 $0.160$ $[0.144,0.174]$예요. 같은 다시 뽑은 표본에서 짝지은 차이는 $+0.092$ $[0.080,0.105]$예요. 등록 분석에서는 라벨을 본 뒤에 본 이 대비가 W1 등록 문서의 가설 F1이 되었어요. \ref{sub:fresh-results}절에서 F1을 고친 점수에 적용해요. 엔도스랭크와 AWP와의 대비도 라벨을 본 뒤에 계산했어요. C-PR은 새 허락 쌍에서 엔도스랭크보다 위에 있고($+0.021$, $[0.013,0.028]$), AWP보다도 위에 있어요($+0.056$, $[0.042,0.070]$). 하지만 새 송금자에서는 AWP보다 한참 뒤져요($-0.085$, $[-0.092,-0.078]$). 세 가지 $\lambda$ 값에 걸쳐 C-PR은 새 허락 쌍에서 거의 평평해요($\lambda=0.25$에서 $0.253$ $[0.243,0.264]$, $\lambda=0.75$에서 $0.249$ $[0.238,0.260]$). 그리고 새 송금자에서 C-PR의 점 추정값은 $\lambda$가 커질수록 $0.269$ $[0.257,0.280]$에서 $0.253$ $[0.242,0.265]$로 내려가요. $\lambda$ 값들 사이의 짝지은 대비는 계산하지 않았고, $\lambda$를 맞추어 고르지도(조정하지도) 않았어요.

씨앗을 쓰는 빼 보기 실험은 허락 신호가 어디에 있어야 하는지 보여 줘요. S-PR은 새 송금자에서 송금 걷기와 거의 같고($-0.001$, $[-0.004,0.002]$), 새 허락 쌍에서는 허락 걷기보다 낮아요($-0.051$, $[-0.066,-0.038]$). S-PR의 다시 출발은 새 허락 쌍에 대한 계수를 송금 걷기의 계수보다 아주 조금만 올려요($0.172$ $[0.157,0.185]$ 대 $0.160$ $[0.144,0.174]$, 짝지은 대비 없음). 그래서 걷기가 어디로 갈지를 송금 화살표가 정하면, 허락 신호는 조금밖에 남지 않아요. 스펜더 홀드아웃에서 허락 화살표는 걷기를 실어 나를 때 도움이 돼요. 허락 걷기에서는 혼자 다 실어 나르고, C-PR에서는 섞여서 실어 날라요. 하지만 걷기가 어디서 다시 출발할지만 정할 때는 거의 도움이 되지 않아요. S-PR의 대비들과, C-PR ($\lambda=1$)과 자기 차수 사이의 차이는 라벨을 본 뒤에 계산했어요. 계획에서 정한 것은 규칙~A의 두 대비뿐이에요.

표 전체에 두 가지 주의할 점이 있어요. 각 라벨은 같은 구성을 미래로 옮긴 것이에요(새 허락, 새 송금자). 그래서 이 스펜더 결과들 가운데 어느 것도 신용이나 손실과는 관계가 없어요. 아래의 트레이더 라벨은 두 종류의 화살표 어느 쪽으로도 만들지 않았지만, 신용 라벨도 아니에요. 또 두 라벨은 모두 듬성듬성해요. 스펜더 17,050개 가운데 2,070개(12.1\%)가 새 허락 쌍을 얻었고, 1,695개(9.9\%)가 새 송금자를 얻었어요. 그래서 모든 계수는 주로 이 작은 무리들이 어떻게 줄 세워지는지에 달려 있어요.

계획에서 정하지 않은 다른 후보 라벨 두 개도 홀드아웃에서 계산했어요. 여기서 보고하지만, 뒷받침하는 증거로 내놓지는 않아요. 허락 취소 수는 모든 점수와 함께 올라가요. 원시 받은 허락 수와 가장 많이 함께 올라가고($0.388$ $[0.376,0.398]$), 그다음은 엔도스랭크예요($0.337$ $[0.327,0.347]$). 송금 위를 걷는 걷기들과는 가장 적게 함께 올라가요. AWP는 $0.191$ $[0.175,0.205]$, C-PR ($\lambda=0$)은 $0.200$ $[0.185,0.214]$예요. 털어 가기 어림 규칙(drain heuristic, 토큰을 빼 가는 일을 대강 찾아내는 규칙)의 값도 모든 점수와 함께 올라가요. 송금 입차수의 $0.112$ $[0.096,0.126]$부터 활동 재시작을 쓴 엔도스랭크의 $0.328$ $[0.318,0.338]$까지예요. 엔도스랭크는 $0.234$ $[0.224,0.245]$, 받은 허락 수는 $0.214$ $[0.200,0.227]$이에요. 보증으로 읽으면 부호가 틀렸어요. 믿을 만한 스펜더라면 허락을 더 자주가 아니라 덜 자주 잃어야 하기 때문이에요. 더 그럴듯한 읽기는 양(얼마나 많이 쓰이는지)이에요. 주인이 많은 스펜더는 주인들이 하는 모든 일을 더 많이 겪어요. 허락 취소와 털어 가기도 마찬가지예요. 그래서 이 라벨들은 믿음과 쓰임을 갈라낼 수 없어요.

\dissertationSubsubheading[sub:trader-results]{트레이더 결과(탐색적)}

\ref{sub:trader-holdout}절의 트레이더 실행은 W0의 동결을 GMX~V2 계약 트레이더(GMX V2에서 거래하는 계약)에 적용해요. 이 실행은 탐색적 결과로 보고해요. 이 실행의 청산 라벨은 \ref{sub:neutral-label-results}절의 비교를 위해 남겨 두었어요. 표~\ref{tab:holdout-traders}에는 4월부터 6월까지 적어도 세 번 포지션을 닫은 계약 트레이더 88개에 대해, 이익 라벨 네 가지에서의 모든 점수가 있어요. 이 가운데 40개는 동결일까지 허락을 받은 적이 있었어요. 트레이더가 88개뿐이라서 모든 구간이 넓어요.

% Generated by experiments/scripts/export_latex.py -- do not edit by hand.
% Source: experiments/data/2-processed-tables-and-evaluations/revised-usd/holdout-w0/eval_summary.json
\begin{table}[htbp]
\centering
\caption[GMX~V2 이익 라벨에 대한 W0 트레이더 분석.]{W0 트레이더 분석이에요. 2026년 3월 31일에 동결한 점수와, 2026년 4월 1일부터 2026년 6월 30일까지 센 GMX~V2 라벨 사이의 켄달의 $\tau$를 보여 줘요. 대상은 동결일 그래프에 있고 그 기간에 세 번 이상 닫은 GMX~V2 계약 계정 $n=88$개예요. 그 가운데 40개는 동결 때까지 허락을 받은 적이 있었어요. 이 분석은 탐색적(미리 정하지 않고 살펴본)이에요. W0의 청산 라벨은 비교를 위해 남겨 두었어요(표~\ref{tab:neutral-label-contrasts}). 이익 본 닫기 수와 실현 이익은 닫은 횟수가 많을수록 커져요. 두 비율은 그렇지 않아요. 신뢰구간(CI)은 계약을 다시 뽑아 구한 95\% 백분위 부트스트랩 구간이에요(짝지은 다시 뽑기 400번, 시드 42).}
\label{tab:holdout-traders}
\fitwidth{%
\begin{tabular}{lcccc}
\toprule
$t_1$의 점수 & 이익 본 닫기 수 & 실현 이익 & 이익 본 닫기 비율 & 손실 없는 닫기 비율 \\
\midrule
$t_1$의 받은 허락 수 (기준선) & \begin{tabular}[c]{@{}c@{}}0.083\\[-1pt]{\footnotesize [-0.103, 0.238]}\end{tabular} & \begin{tabular}[c]{@{}c@{}}0.024\\[-1pt]{\footnotesize [-0.154, 0.202]}\end{tabular} & \begin{tabular}[c]{@{}c@{}}-0.146\\[-1pt]{\footnotesize [-0.317, 0.029]}\end{tabular} & \begin{tabular}[c]{@{}c@{}}-0.278\\[-1pt]{\footnotesize [-0.431, -0.119]}\end{tabular} \\
$t_1$의 송금 입차수 (기준선) & \begin{tabular}[c]{@{}c@{}}0.008\\[-1pt]{\footnotesize [-0.135, 0.169]}\end{tabular} & \begin{tabular}[c]{@{}c@{}}0.008\\[-1pt]{\footnotesize [-0.149, 0.175]}\end{tabular} & \begin{tabular}[c]{@{}c@{}}-0.191\\[-1pt]{\footnotesize [-0.329, -0.041]}\end{tabular} & \begin{tabular}[c]{@{}c@{}}-0.197\\[-1pt]{\footnotesize [-0.350, -0.035]}\end{tabular} \\
\midrule
EndorseRank & \begin{tabular}[c]{@{}c@{}}0.150\\[-1pt]{\footnotesize [-0.019, 0.295]}\end{tabular} & \begin{tabular}[c]{@{}c@{}}0.078\\[-1pt]{\footnotesize [-0.094, 0.226]}\end{tabular} & \begin{tabular}[c]{@{}c@{}}-0.055\\[-1pt]{\footnotesize [-0.196, 0.084]}\end{tabular} & \begin{tabular}[c]{@{}c@{}}-0.178\\[-1pt]{\footnotesize [-0.298, -0.042]}\end{tabular} \\
AWP & \begin{tabular}[c]{@{}c@{}}0.155\\[-1pt]{\footnotesize [0.005, 0.297]}\end{tabular} & \begin{tabular}[c]{@{}c@{}}0.196\\[-1pt]{\footnotesize [0.048, 0.351]}\end{tabular} & \begin{tabular}[c]{@{}c@{}}-0.124\\[-1pt]{\footnotesize [-0.265, 0.015]}\end{tabular} & \begin{tabular}[c]{@{}c@{}}-0.254\\[-1pt]{\footnotesize [-0.404, -0.104]}\end{tabular} \\
\midrule
C-PR ($\lambda=1$) & \begin{tabular}[c]{@{}c@{}}0.083\\[-1pt]{\footnotesize [-0.072, 0.223]}\end{tabular} & \begin{tabular}[c]{@{}c@{}}0.040\\[-1pt]{\footnotesize [-0.133, 0.190]}\end{tabular} & \begin{tabular}[c]{@{}c@{}}-0.074\\[-1pt]{\footnotesize [-0.231, 0.068]}\end{tabular} & \begin{tabular}[c]{@{}c@{}}-0.188\\[-1pt]{\footnotesize [-0.322, -0.047]}\end{tabular} \\
C-PR ($\lambda=0$) & \begin{tabular}[c]{@{}c@{}}0.150\\[-1pt]{\footnotesize [0.011, 0.272]}\end{tabular} & \begin{tabular}[c]{@{}c@{}}0.288\\[-1pt]{\footnotesize [0.128, 0.438]}\end{tabular} & \begin{tabular}[c]{@{}c@{}}-0.101\\[-1pt]{\footnotesize [-0.243, 0.045]}\end{tabular} & \begin{tabular}[c]{@{}c@{}}-0.225\\[-1pt]{\footnotesize [-0.359, -0.083]}\end{tabular} \\
C-PR ($\lambda=0.25$) & \begin{tabular}[c]{@{}c@{}}0.152\\[-1pt]{\footnotesize [0.015, 0.274]}\end{tabular} & \begin{tabular}[c]{@{}c@{}}0.294\\[-1pt]{\footnotesize [0.143, 0.424]}\end{tabular} & \begin{tabular}[c]{@{}c@{}}-0.084\\[-1pt]{\footnotesize [-0.233, 0.069]}\end{tabular} & \begin{tabular}[c]{@{}c@{}}-0.210\\[-1pt]{\footnotesize [-0.352, -0.061]}\end{tabular} \\
C-PR ($\lambda=0.5$) & \begin{tabular}[c]{@{}c@{}}0.160\\[-1pt]{\footnotesize [0.030, 0.284]}\end{tabular} & \begin{tabular}[c]{@{}c@{}}0.283\\[-1pt]{\footnotesize [0.142, 0.410]}\end{tabular} & \begin{tabular}[c]{@{}c@{}}-0.084\\[-1pt]{\footnotesize [-0.227, 0.058]}\end{tabular} & \begin{tabular}[c]{@{}c@{}}-0.208\\[-1pt]{\footnotesize [-0.346, -0.061]}\end{tabular} \\
C-PR ($\lambda=0.75$) & \begin{tabular}[c]{@{}c@{}}0.155\\[-1pt]{\footnotesize [0.025, 0.279]}\end{tabular} & \begin{tabular}[c]{@{}c@{}}0.270\\[-1pt]{\footnotesize [0.127, 0.398]}\end{tabular} & \begin{tabular}[c]{@{}c@{}}-0.082\\[-1pt]{\footnotesize [-0.221, 0.063]}\end{tabular} & \begin{tabular}[c]{@{}c@{}}-0.203\\[-1pt]{\footnotesize [-0.341, -0.059]}\end{tabular} \\
S-PR & \begin{tabular}[c]{@{}c@{}}0.166\\[-1pt]{\footnotesize [0.023, 0.290]}\end{tabular} & \begin{tabular}[c]{@{}c@{}}0.301\\[-1pt]{\footnotesize [0.147, 0.436]}\end{tabular} & \begin{tabular}[c]{@{}c@{}}-0.075\\[-1pt]{\footnotesize [-0.221, 0.072]}\end{tabular} & \begin{tabular}[c]{@{}c@{}}-0.221\\[-1pt]{\footnotesize [-0.365, -0.078]}\end{tabular} \\
\bottomrule
\end{tabular}
}
\end{table}


두 개의 수 라벨(센 수로 된 라벨)은 송금 위를 걷는 걷기들을 다른 점수들과 갈라 놓아요. 이익을 낸 닫기의 수와 실현 이익(실제로 얻은 이익)은 계약이 얼마나 많이 거래하는지에 따라 커져요. 이 두 라벨에서 AWP, C-PR ($\lambda=0$), 세 혼합 점수, S-PR은 구간이 0보다 위에 있는 계수에 이르러요. 이 계수는 이익을 낸 닫기에서 C-PR ($\lambda=0$)의 $0.150$ $[0.011,0.272]$부터, 실현 이익에서 S-PR의 $0.301$ $[0.147,0.436]$까지예요. 반면에 엔도스랭크, C-PR ($\lambda=1$), 두 원시 차수는 구간 안에 0이 들어 있어요. 송금 입차수는 두 라벨 모두에서 0에 가까워요($0.008$ $[-0.135,0.169]$와 $0.008$ $[-0.149,0.175]$). 그러니까 보낸 쪽의 수만으로는 트레이더를 줄 세우지 못해요. 두 비율 라벨은 활동에 따라 커지지 않는데, 여기서는 모든 계수가 음수예요. 손실이 아닌 비율에서는 모든 구간이 0보다 아래에 있어요. 엔도스랭크의 $-0.178$ $[-0.298,-0.042]$부터 받은 허락 수의 $-0.278$ $[-0.431,-0.119]$까지예요. 이익을 낸 비율에서는 송금 입차수의 구간만 0보다 아래에 있어요($-0.191$, $[-0.329,-0.041]$). 이 점수들에서 높게 줄 세워진 계약은, 자기 포지션 가운데 손실 없이 닫은 비율이 더 작았어요. 그리고 그 일부는 손실이 아닌 비율을 만든 방식 때문에 저절로 생겨요(\ref{sub:neutral-label-results}절). 어떤 점수도 계약이 얼마나 거래를 잘하는지는 줄 세우지 못해요. W1에서 규칙~B는 자기만의 트레이더 라벨인 청산 없이 닫은 비율(F3)에서 C-PR을 자기 송금 층과 맞대어 봐요.

\dissertationSubsubheading[sub:fresh-results]{기간 W1: 2026년 7월부터 9월까지}

\ref{sub:fresh-holdout}절의 등록 문서는 등록 분석의 W0 결과를 안 뒤에 썼어요. 그리고 2026년 10월 9일 16:49:32~UTC에 공개 저장소에 커밋했어요(커밋 \texttt{0d3a057}, 커밋의 시간대인 UTC+9로는 10월 10일 01:49). 이 등록 문서는 이 연구가 7월--9월 로그를 뽑기 전에 커밋했어요. \textcite{nosek2018preregistration}에서 말하는 뜻의 사전 등록이에요. 등록 문서는 규칙~B와 그 허용 폭을 정했어요. 또 판정 옆에 함께 보고할 결과 세 개(F4--F6), 민감도 분석 두 개, 그리고 두 종류의 화살표 어느 쪽으로도 만들지 않은 결과 값에 대한 비교도 정했어요. 허용 폭 $\delta=0.01$은 등록 분석에서 F2 대비의 W0 구간의 반폭(구간 너비의 절반)을 소수 둘째 자리까지 올림한 값이에요($0.0040$; $-0.030$, $[-0.034,-0.026]$). 로그를 뽑는 스크립트(추출 스크립트)는 커밋된 등록 문서가 없으면 돌아가지 않아요(부록~\ref{ch:appendix-eval}). 그 매니페스트(실행 기록)에는 등록 커밋, 등록 문서와 설정 파일의 SHA-256 해시가 적혀 있어요. 또 시험용 우회 설정을 쓰지 않았다는 것과, 커밋 1초 뒤인 16:49:33~UTC에 시작했다는 것도 적혀 있어요. 7월--9월 로그에는 계약 계정이 한 GMX~V2 닫기가 6,805개 있었고, 모두 해독되었어요. 아래 결과는 고친 점수의 결과예요. W1의 등록 평가를 안 뒤에, 같은 짝지은 다시 뽑기 400번과 시드 42로 계산했어요. 규칙~B, F4--F6, 두 민감도 분석은 등록한 기준값 그대로 고친 점수에 적용해요. 이것은 재분석이고, 등록한 시험이 아니에요. 등록 분석의 판정은 부록~\ref{sec:map-registered}에 있어요.

2026년 6월 30일 동결일에, 스펜더 계약 18,152개가 다섯 토큰 가운데 하나에 대해 0보다 큰 마지막 허락을 가지고 있었어요. 이 가운데 1,616개가 7월부터 9월까지 새 허락 쌍을 얻었고, 1,533개가 새 송금자를 얻었어요. 동결일 그래프에 있는 GMX~V2 계약 트레이더 70개가 라벨 기간에 적어도 세 번 포지션을 닫았어요. 이 가운데 9개는 청산이 적어도 한 번 있었고, 1개는 모든 닫기가 청산이었어요. 청산이 이렇게 적어서 트레이더 라벨은 별로 달라지지 않고, 트레이더 구간은 모두 넓어요. 표~\ref{tab:fresh-holdout}에는 규칙~B의 모든 점수가 세 라벨에 대해 있고, 표~\ref{tab:fresh-contrasts}에는 대비와 그 판정이 있어요.

% Generated by experiments/scripts/export_latex.py -- do not edit by hand.
% Source: experiments/data/2-processed-tables-and-evaluations/revised-usd/window-w1/eval_summary.json
\begin{table}[htbp]
\centering
\caption[기간 W1: 규칙 B의 세 라벨에 대한 켄달의 $\tau$.]{기간 W1의 결과예요. 2026년 6월 30일에 동결한 점수와 규칙 B의 세 라벨 사이의 켄달의 $\tau$를 보여 줘요. 라벨은 2026년 7월 1일부터 2026년 9월 30일까지 셌어요. 스펜더는 동결 때 마지막 허락이 0보다 큰 계약 스펜더 $n=18{,}152$개예요. 그 가운데 1{,}616개는 새 허락 쌍을 얻었고, 1{,}533개는 새 송금자를 얻었어요. 트레이더는 동결일 그래프에 있고 라벨 기간에 세 번 이상 닫은 GMX~V2 계약 계정 $n=70$개예요. 그 가운데 9개는 청산이 한 번 이상 있었고, 1개는 청산만 있었어요. C-PR ($\lambda=0$)은 규칙 B의 비교 대상이고, AWP는 규칙 B의 첫 번째 민감도 분석의 비교 대상이에요(표~\ref{tab:fresh-contrasts}). 엔도스랭크는 표~\ref{tab:fresh-posthoc}에 있어요. 신뢰구간(CI)은 계약을 다시 뽑아 구한 95\% 백분위 부트스트랩 구간이에요(짝지은 다시 뽑기 400번, 시드 42).}
\label{tab:fresh-holdout}
\fitwidth{%
\begin{tabular}{lccc}
\toprule
 & \multicolumn{2}{c}{스펜더 ($n=18{,}152$)} & 트레이더 ($n=70$) \\
\cmidrule(lr){2-3} \cmidrule(lr){4-4}
2026년 6월 30일의 점수 & 새 허락 쌍 & 새 송금자 & 청산 없이 닫은 비율 \\
\midrule
동결일의 받은 허락 수 (기준선) & 0.268 [0.257, 0.279] & 0.205 [0.193, 0.215] & 0.394 [0.261, 0.524] \\
동결일의 송금 입차수 (기준선) & 0.131 [0.119, 0.143] & 0.308 [0.299, 0.319] & 0.310 [0.185, 0.416] \\
\midrule
AWP & 0.169 [0.157, 0.183] & 0.330 [0.321, 0.340] & 0.329 [0.215, 0.439] \\
\midrule
C-PR ($\lambda=1$) & 0.228 [0.217, 0.238] & 0.188 [0.178, 0.198] & 0.301 [0.177, 0.425] \\
C-PR ($\lambda=0$) & 0.141 [0.129, 0.154] & 0.291 [0.282, 0.301] & 0.350 [0.236, 0.461] \\
C-PR ($\lambda=0.25$) & 0.251 [0.241, 0.260] & 0.270 [0.260, 0.279] & 0.349 [0.234, 0.459] \\
C-PR ($\lambda=0.5$) & 0.250 [0.240, 0.260] & 0.263 [0.254, 0.273] & 0.350 [0.236, 0.461] \\
C-PR ($\lambda=0.75$) & 0.249 [0.238, 0.259] & 0.256 [0.246, 0.266] & 0.347 [0.231, 0.462] \\
S-PR & 0.155 [0.142, 0.168] & 0.291 [0.281, 0.301] & 0.374 [0.255, 0.498] \\
\bottomrule
\end{tabular}
}
\end{table}


% Generated by experiments/scripts/export_latex.py -- do not edit by hand.
% Source: experiments/data/2-processed-tables-and-evaluations/revised-usd/window-w1/eval_summary.json
\begin{table}[htbp]
\centering
\caption[기간 W1: 규칙 B의 대비와 판정, 그리고 민감도 분석.]{기간 W1에서 규칙 B의 짝지은 대비 F1--F6과 그 판정, 그리고 두 가지 민감도 분석을 보여 줘요. 규칙 B와 그 민감도 분석은 원래 단위로 한 분석을 위해 등록했어요. 여기서는 이것들을 고친 점수에 적용했어요. $\Delta\tau=\tau_a-\tau_b$이고, C-PR ($\lambda=0.5$)이 $a$예요. 우월성은 구간이 0보다 위에 있으면 성립해요. 비열등성은 구간의 아래 끝이 $-\delta=-0.01$보다 위에 있으면 성립해요. 규칙 A의 첫 부분인 F6a는 구간이 0을 품거나 0보다 위에 있으면 성립해요. 표~\ref{tab:fresh-holdout}에서처럼 F1, F2, F6a, F6b는 스펜더 계약($n=18{,}152$개)에서, F3--F5는 트레이더($n=70$)에서 구했어요. 신뢰구간(CI)은 계약을 다시 뽑아 구한 95\% 백분위 부트스트랩 구간이에요(짝지은 다시 뽑기 400번, 시드 42).}
\label{tab:fresh-contrasts}
\fitwidth{%
\begin{tabular}{lllcccll}
\toprule
 & 대비 & 라벨 & 검사 & $\tau_a$ & $\tau_b$ & $\Delta\tau$ [95\% 신뢰구간] & 판정 \\
\midrule
\multicolumn{8}{l}{\emph{규칙 B: C-PR ($\lambda=0.5$)과 그 송금 층만 걷는 C-PR ($\lambda=0$)의 비교}} \\
F1 & C-PR 대 C-PR ($\lambda=0$) & 새 허락 쌍 & 우월성 & 0.250 & 0.141 & +0.110 [0.098, 0.122] & 성립 \\
F2 & C-PR 대 C-PR ($\lambda=0$) & 새 송금자 & 비열등성, $-0.01$ & 0.263 & 0.291 & -0.028 [-0.033, -0.023] & 불성립 \\
F3 & C-PR 대 C-PR ($\lambda=0$) & 청산 없이 닫은 비율 & 비열등성, $-0.01$ & 0.350 & 0.350 & +0.000 [-0.014, 0.010] & 불성립 \\
\multicolumn{7}{l}{결정: F1, F2, F3 모두 성립해야 해요} & 충족 안 됨 \\
\midrule
\multicolumn{8}{l}{\emph{보고만 하고 결정에는 들어가지 않음}} \\
F4 & C-PR 대 C-PR ($\lambda=0$) & 청산 없이 닫은 비율 & 우월성 & 0.350 & 0.350 & +0.000 [-0.014, 0.010] & 불성립 \\
F5 & C-PR 대 C-PR ($\lambda=1$) & 청산 없이 닫은 비율 & 우월성 & 0.350 & 0.301 & +0.049 [-0.055, 0.168] & 불성립 \\
F6a & C-PR 대 C-PR ($\lambda=1$) & 새 허락 쌍 & 더 나쁘지 않음 & 0.250 & 0.228 & +0.022 [0.015, 0.029] & 성립 \\
F6b & C-PR 대 C-PR ($\lambda=0$) & 새 송금자 & 우월성 & 0.263 & 0.291 & -0.028 [-0.033, -0.023] & 불성립 \\
\midrule
\multicolumn{8}{l}{\emph{민감도 분석 1: AWP를 비교 대상으로 함}} \\
F1 & C-PR 대 AWP & 새 허락 쌍 & 우월성 & 0.250 & 0.169 & +0.081 [0.069, 0.095] & 성립 \\
F2 & C-PR 대 AWP & 새 송금자 & 비열등성, $-0.01$ & 0.263 & 0.330 & -0.067 [-0.072, -0.061] & 불성립 \\
F3 & C-PR 대 AWP & 청산 없이 닫은 비율 & 비열등성, $-0.01$ & 0.350 & 0.329 & +0.021 [-0.019, 0.060] & 불성립 \\
F4 & C-PR 대 AWP & 청산 없이 닫은 비율 & 우월성 & 0.350 & 0.329 & +0.021 [-0.019, 0.060] & 불성립 \\
\midrule
\multicolumn{8}{l}{\emph{민감도 분석 2: 그래프 밖의 계약을 외톨이 점으로 남김}} \\
F1 & C-PR 대 C-PR ($\lambda=0$) & 새 허락 쌍 & 우월성 & 0.250 & 0.141 & +0.109 [0.098, 0.121] & 성립 \\
F2 & C-PR 대 C-PR ($\lambda=0$) & 새 송금자 & 비열등성, $-0.01$ & 0.263 & 0.291 & -0.028 [-0.033, -0.023] & 불성립 \\
F3 & C-PR 대 C-PR ($\lambda=0$) & 청산 없이 닫은 비율 & 비열등성, $-0.01$ & 0.350 & 0.350 & +0.000 [-0.014, 0.010] & 불성립 \\
F4 & C-PR 대 C-PR ($\lambda=0$) & 청산 없이 닫은 비율 & 우월성 & 0.350 & 0.350 & +0.000 [-0.014, 0.010] & 불성립 \\
F5 & C-PR 대 C-PR ($\lambda=1$) & 청산 없이 닫은 비율 & 우월성 & 0.350 & 0.322 & +0.028 [-0.047, 0.094] & 불성립 \\
F6a & C-PR 대 C-PR ($\lambda=1$) & 새 허락 쌍 & 더 나쁘지 않음 & 0.250 & 0.228 & +0.022 [0.015, 0.029] & 성립 \\
F6b & C-PR 대 C-PR ($\lambda=0$) & 새 송금자 & 우월성 & 0.263 & 0.291 & -0.028 [-0.033, -0.023] & 불성립 \\
\bottomrule
\end{tabular}
}
\end{table}


고친 점수에 적용하면 규칙~B는 충족되지 않아요. F1은 성립해요. C-PR은 나중에 새 허락 쌍을 얻은 스펜더를 $\tau=0.250$ $[0.240,0.260]$으로 줄 세우고, 자기 송금 층만 걷는 걷기는 $0.141$ $[0.129,0.154]$예요. 차이는 $+0.110$ $[0.098,0.122]$예요. 그러니까 W0에서 얻은 $+0.092$ $[0.080,0.105]$의 이득이 W1에서도 다시 나타나요. F2는 성립하지 않아요. 새 송금자에서 C-PR은 $0.263$ $[0.254,0.273]$이고, 송금 걷기는 $0.291$ $[0.282,0.301]$이에요. 차이는 $-0.028$ $[-0.033,-0.023]$이고, 구간 전체가 허용 폭 $-0.01$보다 아래에 있어요. 그러니까 C-PR은 나중에 들어오는 흐름을 송금 걷기보다 못하게 줄 세우고, 그 차이는 등록 문서가 허용한 것보다 커요. F3은 아슬아슬하게 성립하지 않아요. 트레이더의 청산 없이 닫은 비율에서 C-PR과 송금 걷기는 둘 다 $0.350$ $[0.236,0.461]$에 이르러요. 차이는 $+0.000$ $[-0.014,0.010]$이고, 그 아래 끝이 $-0.01$보다 아래에 있어요. 규칙은 세 가설이 모두 성립해야 했어요. 그래서 등록 문서가 이런 경우에 대해 말한 대로, 이 논문은 C-PR이 송금만 걷는 걷기보다 나아진다고 주장하지 않아요. 등록 분석도 같은 판정에 이르렀어요(표~\ref{tab:registered-vs-revised}).

등록 문서는 두 계수가 모두 0 가까이에 있어도 F3이 통과할 수 있다고 경고했어요. 여기서는 어느 계수도 0 가까이에 있지 않아요. F3이 실패한 까닭은, 점 추정값이 같은데도 짝지은 구간이 허용 폭 아래까지 내려가기 때문이에요. 모든 점수가 이 라벨을 0보다 위에 있는 구간으로 줄 세워요. 동결일의 받은 허락 수는 $0.394$ $[0.261,0.524]$, S-PR은 $0.374$ $[0.255,0.498]$에 이르러요. C-PR의 세 혼합 점수는 $0.347$ $[0.231,0.462]$와 $0.350$ $[0.236,0.461]$ 사이에 있어요. AWP는 $0.329$ $[0.215,0.439]$, 송금 입차수는 $0.310$ $[0.185,0.416]$, C-PR ($\lambda=1$)은 $0.301$ $[0.177,0.425]$예요(표~\ref{tab:fresh-holdout}).

판정 옆에 함께 보고하는 결과는 등록한 순서를 따라요. F4는 성립하지 않아요. 청산 없이 닫은 비율에서 C-PR은 송금 걷기보다 낫지 않고, 차이의 구간 안에 0이 들어 있어요. F5도 성립하지 않아요. 그 라벨에서 C-PR과 허락 걷기를 서로 가려낼 수 없기 때문이에요($+0.049$, $[-0.055,0.168]$). F6은 규칙~A의 주된 부분을 W1에서 다시 해 본 것인데, W0에서처럼 성립하지 않아요. C-PR은 새 허락 쌍에서 허락 걷기보다 나쁘지 않지만($+0.022$, $[0.015,0.029]$), 새 송금자에서는 송금 걷기보다 아래에 있어요($-0.028$, $[-0.033,-0.023]$). W0에서처럼 C-PR의 점 추정값은 새 허락 쌍에서는 $\lambda$에 걸쳐 거의 평평하고, 새 송금자에서는 $\lambda$가 커질수록 내려가요(표~\ref{tab:fresh-holdout}). S-PR은 새 허락 쌍을 $0.155$ $[0.142,0.168]$로 줄 세워, 허락 걷기($0.228$ $[0.217,0.238]$)보다 낮아요. 그리고 새 송금자를 $0.291$ $[0.281,0.301]$로 줄 세워, 송금 걷기($0.291$ $[0.282,0.301]$)와 같은 수준이에요. W1에서 S-PR의 짝지은 대비는 하나도 등록하지 않았어요.

두 민감도 분석 어느 쪽도 판정을 바꾸지 않아요(표~\ref{tab:fresh-contrasts}). AWP를 비교 대상으로 삼으면, F1은 성립하고($+0.081$, $[0.069,0.095]$) F2는 뚜렷하게 성립하지 않아요. AWP는 새 송금자를 $0.330$ $[0.321,0.340]$으로 줄 세우고, C-PR과 AWP의 차이는 $-0.067$ $[-0.072,-0.061]$이에요. F3과 F4도 성립하지 않아요($+0.021$, $[-0.019,0.060]$). 그러니까 AWP와 견주면, C-PR은 나중의 허락에서 얻는 이득과 맞바꾸어, 나중에 들어오는 흐름을 줄 세우는 힘 가운데 눈에 띄는 만큼을 내줘요. 그래프 밖에 있는 계약을 외톨이 점으로 남겨 두면, 모든 판정이 그대로예요. 그리고 F1부터 F4까지와 F6의 점 추정값은 많아야 $0.001$만큼 바뀌어요. F5는 여전히 가려지지 않아요. 이 규칙은 허락 그래프 밖에 있는 트레이더의 점수를 0에서, 들어오는 화살표가 없는 점의 점수로 올려 줘요. 그러면 청산 없이 닫은 비율에 대한 허락 걷기의 계수가 올라가고, C-PR이 허락 걷기보다 앞선 정도는 $+0.028$ $[-0.047,0.094]$로 줄어요.

표~\ref{tab:fresh-posthoc}에는 엔도스랭크와, 활동 재시작을 쓴 그 변형이 더해져 있어요. 둘 다 규칙~B의 점수가 아니고, 그 행들은 판정에 하나도 들어가지 않아요. 이 표의 대비 가운데 어느 것도 등록 문서에서 정하지 않았어요. 모두 7월--9월 라벨을 본 뒤에 계산했고, 변형의 대비는 그보다 더 뒤에 저장해 둔 점수로 계산했어요. 한 층만 걷는 두 걷기 모두에서 W0의 모양이 그대로 나타나요. 동결일의 받은 허락 수는 이번에도 새 허락 쌍을 엔도스랭크보다 더 잘 줄 세워요($0.268$ $[0.257,0.279]$ 대 $0.234$ $[0.222,0.244]$; $-0.034$, $[-0.038,-0.029]$). 그리고 AWP는 이번에도 새 송금자를 송금 입차수보다 더 잘 줄 세워요($0.330$ $[0.321,0.340]$ 대 $0.308$ $[0.299,0.319]$; $+0.022$, $[0.017,0.027]$). 활동 재시작을 쓴 엔도스랭크는 W0에서처럼 새 허락 쌍을 차수보다 더 잘 줄 세워요($+0.027$, $[0.021,0.034]$). C-PR은 W0에서처럼 새 허락 쌍을 엔도스랭크보다 더 잘 줄 세워요($+0.016$, $[0.010,0.024]$). 그리고 청산 없이 닫은 비율에서는 둘을 서로 가려낼 수 없어요($+0.053$, $[-0.051,0.171]$).

% Generated by experiments/scripts/export_latex.py -- do not edit by hand.
% Source: experiments/data/2-processed-tables-and-evaluations/revised-usd/window-w1/eval_summary.json
% Source: experiments/data/2-processed-tables-and-evaluations/revised-usd/supplementary/supplementary.json
\begin{table}[htbp]
\centering
\caption[기간 W1에서 엔도스랭크와 활동 재시작을 쓰는 그 변형.]{기간 W1의 코호트와 라벨에서 엔도스랭크와, 활동 재시작을 쓰는 엔도스랭크 변형을 보여 줘요(스펜더 계약 $n=18{,}152$개, 트레이더 $n=70$). 이 행들도, 그 아래의 대비들도 규칙 B에 들지 않아요. 그리고 이 대비들은 등록 문서에 이름이 적혀 있지 않았어요. 활동 재시작을 쓰는 엔도스랭크의 대비는 저장해 둔 W1 점수로 나중에 계산했어요. AWP, C-PR, 차수 행은 표~\ref{tab:fresh-holdout}의 값을 다시 적은 것이고, 모든 대비는 같은 계약 다시 뽑은 표본을 써요. 엔도스랭크와 AWP 사이의 차이는 표~\ref{tab:endorserank-awp-holdout}에 있어요. 신뢰구간(CI)은 계약을 다시 뽑아 구한 95\% 백분위 부트스트랩 구간이에요(짝지은 다시 뽑기 400번, 시드 42).}
\label{tab:fresh-posthoc}
\fitwidth{%
\begin{tabular}{lccc}
\toprule
 & \multicolumn{2}{c}{스펜더 ($n=18{,}152$)} & 트레이더 ($n=70$) \\
\cmidrule(lr){2-3} \cmidrule(lr){4-4}
2026년 6월 30일의 점수 & 새 허락 쌍 & 새 송금자 & 청산 없이 닫은 비율 \\
\midrule
동결일의 받은 허락 수 (기준선) & 0.268 [0.257, 0.279] & 0.205 [0.193, 0.215] & 0.394 [0.261, 0.524] \\
동결일의 송금 입차수 (기준선) & 0.131 [0.119, 0.143] & 0.308 [0.299, 0.319] & 0.310 [0.185, 0.416] \\
\midrule
EndorseRank & 0.234 [0.222, 0.244] & 0.194 [0.184, 0.205] & 0.297 [0.174, 0.418] \\
EndorseRank, 활동 재시작(허락하는 활동) & 0.295 [0.285, 0.304] & 0.242 [0.232, 0.253] & 0.288 [0.161, 0.405] \\
AWP & 0.169 [0.157, 0.183] & 0.330 [0.321, 0.340] & 0.329 [0.215, 0.439] \\
C-PR ($\lambda=0.5$) & 0.250 [0.240, 0.260] & 0.263 [0.254, 0.273] & 0.350 [0.236, 0.461] \\
\midrule
EndorseRank 빼기 받은 허락 수, 새 허락 쌍 & \multicolumn{3}{l}{$\Delta\tau=-0.034$ [-0.038, -0.029]} \\
AWP 빼기 송금 입차수, 새 송금자 & \multicolumn{3}{l}{$\Delta\tau=+0.022$ [0.017, 0.027]} \\
C-PR 빼기 EndorseRank, 새 허락 쌍 & \multicolumn{3}{l}{$\Delta\tau=+0.016$ [0.010, 0.024]} \\
C-PR 빼기 EndorseRank, 청산 없이 닫은 비율 & \multicolumn{3}{l}{$\Delta\tau=+0.053$ [-0.051, 0.171]} \\
\midrule
EndorseRank (활동 재시작) 빼기 받은 허락 수, 새 허락 쌍 & \multicolumn{3}{l}{$\Delta\tau=+0.027$ [0.021, 0.034]} \\
EndorseRank (활동 재시작) 빼기 EndorseRank, 새 허락 쌍 & \multicolumn{3}{l}{$\Delta\tau=+0.061$ [0.054, 0.067]} \\
EndorseRank (활동 재시작) 빼기 EndorseRank, 새 송금자 & \multicolumn{3}{l}{$\Delta\tau=+0.048$ [0.042, 0.053]} \\
\bottomrule
\end{tabular}
}
\end{table}


표~\ref{tab:fresh-traders}에는 판정에 들어가지 않는 W1의 트레이더 라벨이 있어요. \ref{sub:neutral-label-results}절의 다시 뽑기 2,000번으로 계산했어요. 송금을 바탕으로 만든 점수들이 두 수 라벨을 줄 세워요. 이익을 낸 닫기에서 AWP는 $0.288$ $[0.125,0.442]$, C-PR ($\lambda=0$)은 $0.199$ $[0.043,0.355]$에 이르러요. 실현 이익에서 AWP는 $0.346$ $[0.189,0.490]$, C-PR ($\lambda=0$)은 $0.338$ $[0.178,0.485]$에 이르러요. 반면에 허락 쪽 점수들은 두 라벨 모두에서 0 가까이에 머물러요. 청산이 없었는지 나타내는 표시에서는 모든 점수가 양수예요. 받은 허락 수가 $0.407$ $[0.268,0.541]$로 가장 높고, 활동 재시작을 쓴 엔도스랭크가 $0.294$ $[0.158,0.417]$로 가장 낮아요. 두 비율 라벨에서는 허락 쪽 점수들이 음수예요. AWP와 C-PR ($\lambda=0$)은 양수이지만, 구간 안에 0이 들어 있어요. 손실이 아닌 비율에서는 네 점수의 구간이 0보다 아래에 있어요. 받은 허락 수는 $-0.294$, $[-0.468,-0.106]$, C-PR ($\lambda=1$)은 $-0.287$, $[-0.443,-0.121]$, 엔도스랭크는 $-0.264$, $[-0.408,-0.106]$, 활동 재시작을 쓴 엔도스랭크 변형은 $-0.269$, $[-0.412,-0.118]$이에요. 이익을 낸 비율에서도 같은 네 점수의 구간이 0보다 아래에 있어요($-0.221$, $[-0.406,-0.034]$; $-0.208$, $[-0.370,-0.030]$; $-0.180$, $[-0.340,-0.012]$; $-0.199$, $[-0.346,-0.039]$).

% Generated by experiments/scripts/export_latex.py -- do not edit by hand.
% Source: experiments/data/2-processed-tables-and-evaluations/revised-usd/window-w1/eval_summary.json
\begin{table}[htbp]
\centering
\caption[기간 W1: 다른 트레이더 라벨에 대한 켄달의 $\tau$.]{기간 W1의 트레이더예요. 2026년 6월 30일에 동결한 점수와, 2026년 7월 1일부터 2026년 9월 30일까지 센 다른 GMX~V2 라벨 사이의 켄달의 $\tau$를 보여 줘요. 대상은 그 기간에 포지션을 세 번 이상 닫은 계약 계정 $n=70$개예요. “청산 없음”은 그 계정이 닫은 포지션 가운데 청산이 하나도 없을 때 1이에요. 이익 본 닫기 수와 실현 이익은 닫은 횟수가 많을수록 커져요. 두 비율은 그렇지 않아요. 이 라벨들은 보고만 하고 규칙 B에는 들어가지 않아요. 구간은 표~\ref{tab:neutral-label-contrasts}에 쓴 다시 뽑은 표본을 써요. 신뢰구간(CI)은 계약을 다시 뽑아 구한 95\% 백분위 부트스트랩 구간이에요(짝지은 다시 뽑기 2{,}000번, 시드 42).}
\label{tab:fresh-traders}
\fitwidth{%
\begin{tabular}{lccccc}
\toprule
2026년 6월 30일의 점수 & 청산 없음 & 이익 본 닫기 수 & 실현 이익 & 이익 본 닫기 비율 & 손실 없는 닫기 비율 \\
\midrule
동결일의 받은 허락 수 (기준선) & \begin{tabular}[c]{@{}c@{}}0.407\\[-1pt]{\footnotesize [0.268, 0.541]}\end{tabular} & \begin{tabular}[c]{@{}c@{}}0.042\\[-1pt]{\footnotesize [-0.158, 0.243]}\end{tabular} & \begin{tabular}[c]{@{}c@{}}0.033\\[-1pt]{\footnotesize [-0.164, 0.230]}\end{tabular} & \begin{tabular}[c]{@{}c@{}}-0.221\\[-1pt]{\footnotesize [-0.406, -0.034]}\end{tabular} & \begin{tabular}[c]{@{}c@{}}-0.294\\[-1pt]{\footnotesize [-0.468, -0.106]}\end{tabular} \\
동결일의 송금 입차수 (기준선) & \begin{tabular}[c]{@{}c@{}}0.329\\[-1pt]{\footnotesize [0.189, 0.451]}\end{tabular} & \begin{tabular}[c]{@{}c@{}}0.185\\[-1pt]{\footnotesize [0.014, 0.341]}\end{tabular} & \begin{tabular}[c]{@{}c@{}}0.168\\[-1pt]{\footnotesize [0.028, 0.291]}\end{tabular} & \begin{tabular}[c]{@{}c@{}}-0.025\\[-1pt]{\footnotesize [-0.194, 0.131]}\end{tabular} & \begin{tabular}[c]{@{}c@{}}-0.042\\[-1pt]{\footnotesize [-0.203, 0.114]}\end{tabular} \\
\midrule
EndorseRank & \begin{tabular}[c]{@{}c@{}}0.303\\[-1pt]{\footnotesize [0.172, 0.422]}\end{tabular} & \begin{tabular}[c]{@{}c@{}}0.073\\[-1pt]{\footnotesize [-0.115, 0.260]}\end{tabular} & \begin{tabular}[c]{@{}c@{}}0.000\\[-1pt]{\footnotesize [-0.198, 0.200]}\end{tabular} & \begin{tabular}[c]{@{}c@{}}-0.180\\[-1pt]{\footnotesize [-0.340, -0.012]}\end{tabular} & \begin{tabular}[c]{@{}c@{}}-0.264\\[-1pt]{\footnotesize [-0.408, -0.106]}\end{tabular} \\
EndorseRank, 허락하는 활동에 따라 다시 출발 & \begin{tabular}[c]{@{}c@{}}0.294\\[-1pt]{\footnotesize [0.158, 0.417]}\end{tabular} & \begin{tabular}[c]{@{}c@{}}0.057\\[-1pt]{\footnotesize [-0.123, 0.239]}\end{tabular} & \begin{tabular}[c]{@{}c@{}}-0.004\\[-1pt]{\footnotesize [-0.179, 0.181]}\end{tabular} & \begin{tabular}[c]{@{}c@{}}-0.199\\[-1pt]{\footnotesize [-0.346, -0.039]}\end{tabular} & \begin{tabular}[c]{@{}c@{}}-0.269\\[-1pt]{\footnotesize [-0.412, -0.118]}\end{tabular} \\
AWP & \begin{tabular}[c]{@{}c@{}}0.333\\[-1pt]{\footnotesize [0.207, 0.447]}\end{tabular} & \begin{tabular}[c]{@{}c@{}}0.288\\[-1pt]{\footnotesize [0.125, 0.442]}\end{tabular} & \begin{tabular}[c]{@{}c@{}}0.346\\[-1pt]{\footnotesize [0.189, 0.490]}\end{tabular} & \begin{tabular}[c]{@{}c@{}}0.136\\[-1pt]{\footnotesize [-0.039, 0.300]}\end{tabular} & \begin{tabular}[c]{@{}c@{}}0.082\\[-1pt]{\footnotesize [-0.077, 0.247]}\end{tabular} \\
\midrule
C-PR ($\lambda=1$) & \begin{tabular}[c]{@{}c@{}}0.307\\[-1pt]{\footnotesize [0.174, 0.428]}\end{tabular} & \begin{tabular}[c]{@{}c@{}}0.010\\[-1pt]{\footnotesize [-0.167, 0.192]}\end{tabular} & \begin{tabular}[c]{@{}c@{}}-0.066\\[-1pt]{\footnotesize [-0.249, 0.121]}\end{tabular} & \begin{tabular}[c]{@{}c@{}}-0.208\\[-1pt]{\footnotesize [-0.370, -0.030]}\end{tabular} & \begin{tabular}[c]{@{}c@{}}-0.287\\[-1pt]{\footnotesize [-0.443, -0.121]}\end{tabular} \\
C-PR ($\lambda=0$) & \begin{tabular}[c]{@{}c@{}}0.358\\[-1pt]{\footnotesize [0.228, 0.475]}\end{tabular} & \begin{tabular}[c]{@{}c@{}}0.199\\[-1pt]{\footnotesize [0.043, 0.355]}\end{tabular} & \begin{tabular}[c]{@{}c@{}}0.338\\[-1pt]{\footnotesize [0.178, 0.485]}\end{tabular} & \begin{tabular}[c]{@{}c@{}}0.136\\[-1pt]{\footnotesize [-0.034, 0.290]}\end{tabular} & \begin{tabular}[c]{@{}c@{}}0.038\\[-1pt]{\footnotesize [-0.127, 0.201]}\end{tabular} \\
C-PR ($\lambda=0.25$) & \begin{tabular}[c]{@{}c@{}}0.357\\[-1pt]{\footnotesize [0.227, 0.474]}\end{tabular} & \begin{tabular}[c]{@{}c@{}}0.155\\[-1pt]{\footnotesize [0.002, 0.305]}\end{tabular} & \begin{tabular}[c]{@{}c@{}}0.284\\[-1pt]{\footnotesize [0.122, 0.435]}\end{tabular} & \begin{tabular}[c]{@{}c@{}}0.090\\[-1pt]{\footnotesize [-0.075, 0.242]}\end{tabular} & \begin{tabular}[c]{@{}c@{}}-0.011\\[-1pt]{\footnotesize [-0.173, 0.148]}\end{tabular} \\
C-PR ($\lambda=0.5$) & \begin{tabular}[c]{@{}c@{}}0.358\\[-1pt]{\footnotesize [0.229, 0.475]}\end{tabular} & \begin{tabular}[c]{@{}c@{}}0.114\\[-1pt]{\footnotesize [-0.044, 0.267]}\end{tabular} & \begin{tabular}[c]{@{}c@{}}0.240\\[-1pt]{\footnotesize [0.075, 0.397]}\end{tabular} & \begin{tabular}[c]{@{}c@{}}0.053\\[-1pt]{\footnotesize [-0.115, 0.212]}\end{tabular} & \begin{tabular}[c]{@{}c@{}}-0.053\\[-1pt]{\footnotesize [-0.217, 0.103]}\end{tabular} \\
C-PR ($\lambda=0.75$) & \begin{tabular}[c]{@{}c@{}}0.353\\[-1pt]{\footnotesize [0.226, 0.472]}\end{tabular} & \begin{tabular}[c]{@{}c@{}}0.081\\[-1pt]{\footnotesize [-0.078, 0.244]}\end{tabular} & \begin{tabular}[c]{@{}c@{}}0.201\\[-1pt]{\footnotesize [0.023, 0.368]}\end{tabular} & \begin{tabular}[c]{@{}c@{}}0.016\\[-1pt]{\footnotesize [-0.153, 0.182]}\end{tabular} & \begin{tabular}[c]{@{}c@{}}-0.090\\[-1pt]{\footnotesize [-0.256, 0.072]}\end{tabular} \\
S-PR & \begin{tabular}[c]{@{}c@{}}0.381\\[-1pt]{\footnotesize [0.248, 0.495]}\end{tabular} & \begin{tabular}[c]{@{}c@{}}0.167\\[-1pt]{\footnotesize [0.018, 0.312]}\end{tabular} & \begin{tabular}[c]{@{}c@{}}0.286\\[-1pt]{\footnotesize [0.122, 0.441]}\end{tabular} & \begin{tabular}[c]{@{}c@{}}0.100\\[-1pt]{\footnotesize [-0.048, 0.244]}\end{tabular} & \begin{tabular}[c]{@{}c@{}}-0.011\\[-1pt]{\footnotesize [-0.164, 0.151]}\end{tabular} \\
\bottomrule
\end{tabular}
}
\end{table}


\dissertationSubsubheading[sub:neutral-label-results]{두 종류의 화살표 어느 쪽으로도 만들지 않은 결과 값}

GMX 트레이더 라벨은 이 자료에서 두 종류의 화살표 어느 쪽으로도 만들지 않은 유일한 결과 값이에요. \ref{sub:neutral-label-protocol}절의 비교는 규칙~B와 함께 등록했어요. 그때는 이 연구가 7월--9월 로그를 뽑기 전이었고, 어떤 점수에 대해서도 W0의 청산 라벨을 계산하기 전이었어요. 여기서는 등록 분석의 결과를 안 뒤에, 그 비교의 대비, 규칙, 문구를 고친 점수에 적용해요. 아래의 판정은 재분석이고, 등록 분석도 같은 판정에 이르렀어요(부록~\ref{sec:map-registered}). W1이 주된 기간이고, W0이 보조 기간이에요. 구간은 짝지은 다시 뽑기 2,000번에서 나왔어요. 97.5\%라고 표시하지 않은 구간은 95\% 구간이에요. 97.5\% 구간은 판정 대비 P와 Q의 본페로니 구간(여러 번 비교할 때 넓힌 신뢰구간)이에요. 표~\ref{tab:neutral-label-contrasts}에는 청산 없이 닫은 비율에서의 대비가 있어요. 표~\ref{tab:neutral-label-recipients}에는 허락을 받는 트레이더가 있고, 표~\ref{tab:neutral-label-grid}(부록~\ref{sec:er-awp-neutral-grid})에는 트레이더 라벨 여섯 가지가 모두 있어요. 표본이 작아요. W1에서는 계약 트레이더 70개, W0에서는 88개예요. 그래서 아래의 모든 구간이 넓어요.

% Generated by experiments/scripts/export_latex.py -- do not edit by hand.
% Source: experiments/data/2-processed-tables-and-evaluations/revised-usd/window-w1/eval_summary.json
\begin{table}[htbp]
\centering
\caption[청산 없이 닫은 비율에서 허락 쪽 점수와 송금 쪽 점수의 비교.]{허락 쪽 점수와 송금 쪽 점수를 청산 없이 닫은 비율에서 서로 비교했어요. 대상은 라벨 기간에 포지션을 세 번 이상 닫은 GMX~V2 계약 트레이더예요. 이 라벨은 두 화살표 종류 어느 것으로도 만들지 않았어요(W1의 트레이더 $n=70$, W0의 트레이더 $n=88$). 비교 쌍, 판정 대비 P, Q와 그 본페로니 97.5\% 구간, 닫은 횟수로 묶음을 나눈(층화한) 행은 이 연구가 7월부터 9월까지의 로그를 뽑기 전에 등록했어요. 이것들은 원래 단위로 한 분석을 위해 등록했고, 여기서는 고친 점수에 적용했어요. W1이 주된 기간이고, W0이 보조 기간이에요. 계약을 단위로 한 짝지은 부트스트랩이고, 다시 뽑기는 2{,}000번이에요(시드 42). 층화한 행에서는, 라벨 기간에 닫은 횟수(3--4, 5--9, 10--24, 25 이상)에 따라 묶음을 나누고 묶음마다 켄달의 $\tau_b$를 구해요. 그리고 계약 쌍의 수를 무게로 삼아 이 값들을 합쳐요. 등록한 규칙을 고친 점수에 적용한 결과는 이래요. R1(그냥 센 수의 수준에서 앞섬)은 충족되지 않았어요. R2(페이지랭크 수준의 앞섬)도 충족되지 않았어요. R3(특이성) 규칙은 적용돼요.}
\label{tab:neutral-label-contrasts}
\fitwidth{%
\begin{tabular}{lccccc}
\toprule
대비 ($a$ 빼기 $b$) & $\tau_a$ & $\tau_b$ & $\Delta\tau$ & 95\% 신뢰구간 & 97.5\% 신뢰구간 \\
\midrule
\multicolumn{6}{l}{\emph{기간 W1, 2026년 7월--9월 (트레이더 $n=70$)}} \\
받은 허락 수 빼기 송금 입차수 & 0.394 & 0.310 & 0.084 & [-0.013, 0.176] & [-0.026, 0.188] \\
EndorseRank (활동 재시작) 빼기 AWP & 0.288 & 0.329 & -0.041 & [-0.135, 0.050] & [-0.148, 0.064] \\
C-PR ($\lambda=1$) 빼기 C-PR ($\lambda=0$) & 0.301 & 0.350 & -0.049 & [-0.169, 0.060] & -- \\
EndorseRank 빼기 AWP & 0.297 & 0.329 & -0.032 & [-0.133, 0.066] & -- \\
C-PR ($\lambda=0$) 빼기 AWP & 0.350 & 0.329 & 0.021 & [-0.016, 0.065] & -- \\
EndorseRank 빼기 받은 허락 수 & 0.297 & 0.394 & -0.097 & [-0.174, -0.009] & -- \\
\quad 받은 허락 수 빼기 송금 입차수, 닫은 횟수로 층화함 & 0.381 & 0.280 & 0.101 & [-0.024, 0.231] & -- \\
\quad EndorseRank (활동 재시작) 빼기 AWP, 닫은 횟수로 층화함 & 0.273 & 0.301 & -0.028 & [-0.207, 0.098] & -- \\
\midrule
\multicolumn{6}{l}{\emph{기간 W0, 2026년 4월--6월 (트레이더 $n=88$)}} \\
받은 허락 수 빼기 송금 입차수 & 0.344 & 0.161 & 0.182 & [0.052, 0.314] & [0.032, 0.338] \\
EndorseRank (활동 재시작) 빼기 AWP & 0.227 & 0.240 & -0.014 & [-0.147, 0.130] & [-0.161, 0.151] \\
C-PR ($\lambda=1$) 빼기 C-PR ($\lambda=0$) & 0.225 & 0.265 & -0.041 & [-0.187, 0.105] & -- \\
EndorseRank 빼기 AWP & 0.224 & 0.240 & -0.017 & [-0.156, 0.128] & -- \\
C-PR ($\lambda=0$) 빼기 AWP & 0.265 & 0.240 & 0.025 & [-0.022, 0.077] & -- \\
EndorseRank 빼기 받은 허락 수 & 0.224 & 0.344 & -0.120 & [-0.203, -0.033] & -- \\
\quad 받은 허락 수 빼기 송금 입차수, 닫은 횟수로 층화함 & 0.338 & 0.209 & 0.129 & [0.020, 0.293] & -- \\
\quad EndorseRank (활동 재시작) 빼기 AWP, 닫은 횟수로 층화함 & 0.215 & 0.252 & -0.037 & [-0.173, 0.125] & -- \\
\bottomrule
\end{tabular}
}
\end{table}


R1은 성립하지 않아요. W1에서 동결일의 받은 허락 수는 청산 없이 닫은 비율에서 $\tau=0.394$ $[0.265,0.518]$에 이르고, 송금 입차수는 $0.310$ $[0.184,0.425]$에 이르러요. 그래서 둘의 차이인 P는 $+0.084$ $[-0.013,0.176]$이고, 97.5\% 구간은 $[-0.026,0.188]$로 0을 포함해요. 닫기 횟수로 나눈 묶음 안에서 계산하면(층화하면), P는 $+0.101$ $[-0.024,0.231]$이에요. W0에서 P는 $+0.182$(97.5\% 구간 $[0.032,0.338]$)로, R1이 요구한 대로 양수예요. 하지만 이 규칙은 W1의 두 구간도 모두 0보다 위에 있어야 했어요. 등록한 문구로 말하면, 나중에 청산을 피하는 데에서 그냥 센 수의 수준으로는 허락이 송금보다 앞서지 않아요. R2도 성립하지 않아요. Q는 AWP의 활동 재시작을 쓴 엔도스랭크에서 AWP를 뺀 값이에요(조건을 똑같이 맞춘 페이지랭크 짝). Q는 W1에서 $-0.041$(97.5\% 구간 $[-0.148,0.064]$), W0에서 $-0.014$(97.5\% 구간 $[-0.161,0.151]$)예요. 그래서 등록한 문구로 말하면, 페이지랭크 버전들은 어느 쪽이 나은지 가려지지 않아요. 다른 짝들도 가려지지 않아요. C-PR의 두 층은 W1에서 $-0.049$ $[-0.169,0.060]$, W0에서 $-0.041$ $[-0.187,0.105]$만큼 차이 나요. 엔도스랭크와 AWP는 $-0.032$ $[-0.133,0.066]$와 $-0.017$ $[-0.156,0.128]$만큼, C-PR ($\lambda=0$)과 AWP는 $+0.021$ $[-0.016,0.065]$와 $+0.025$ $[-0.022,0.077]$만큼 차이 나요. 등록 문서가 정하지 않았고 라벨을 본 뒤에 계산한 대비에서는, 엔도스랭크가 두 기간 모두에서 허락한 주인 수(서로 다른 주인 수)보다 아래에 있어요. W0에서는 $-0.120$ $[-0.203,-0.033]$, W1에서는 $-0.097$ $[-0.174,-0.009]$만큼이에요.

R3은 적용돼요. W1에서 P는 이익을 낸 비율($-0.196$, $[-0.364,-0.046]$)과 손실이 아닌 비율($-0.252$, $[-0.391,-0.120]$)에서 0보다 아래에 있어요. 그래서 등록한 문구로 말하면, 이 신호는 거래를 잘하는 것이 아니라 강제 청산을 피하는 것에 관한 거예요. 손실이 아닌 비율의 차이 가운데 일부는 만든 방식 때문에 저절로 생겨요. 처리 과정은 청산된 닫기를 손실이 아니었던 닫기로 세요. 그래서 청산이 있는 트레이더는 손실이 아닌 비율이 올라가요. 반면에 이익을 낸 비율은 청산된 닫기를 이익을 내지 못한 닫기로 세요. 그래서 R3은 주로 이익을 낸 비율에 바탕을 두어요. 다른 라벨들은 판정에 들어가지 않아요(표~\ref{tab:neutral-label-grid}). 청산이 없었는지 나타내는 표시에서 P는 W1에서 $+0.078$ $[-0.012,0.166]$, W0에서 $+0.197$ $[0.059,0.333]$이에요. 활동에 따라 커지는 이익을 낸 닫기와 실현 이익에서는, W1에서 페이지랭크 짝들이 송금 쪽을 편들어요. 구간이 0보다 아래에 있는 것은 Q($-0.231$, $[-0.429,-0.027]$과 $-0.350$, $[-0.536,-0.157]$), 엔도스랭크 대 AWP($-0.214$, $[-0.423,-0.012]$와 $-0.346$, $[-0.547,-0.144]$), 그리고 실현 이익에서 C-PR의 두 층($-0.404$, $[-0.624,-0.182]$)이에요. W1에서는 같은 세 짝이 두 비율 라벨에서도 송금 쪽을 편들어요(표~\ref{tab:neutral-label-grid}). P는 두 수 라벨에서 같은 쪽을 가리키지만, 가려지지는 않아요. 그리고 W0에서는 허락 쪽 점수와 송금 쪽 점수의 짝 가운데, 실현 이익에서 C-PR의 두 층만 가려져요($-0.248$, $[-0.433,-0.067]$).

% Generated by experiments/scripts/export_latex.py -- do not edit by hand.
% Source: experiments/data/2-processed-tables-and-evaluations/revised-usd/window-w1/eval_summary.json
\begin{table}[htbp]
\centering
\caption[허락을 받는 트레이더와 라벨이 있는 다른 모든 트레이더의 비교.]{동결일에 받은 허락 수가 0보다 큰 GMX~V2 계약 트레이더(받는 쪽)를, 라벨이 있는 다른 모든 트레이더와 두 기간에서 비교했어요(W1의 트레이더 $n=70$, W0의 트레이더 $n=88$). 받는 쪽과 다른 트레이더의 비교, 그리고 영향 점검은 표~\ref{tab:neutral-label-contrasts}의 대비와 함께 등록했어요. 여기서는 고친 점수에 적용했어요. 트레이더 가운데 가장 큰 받은 허락 수는 W1에서 1, W0에서 3이에요. 그래서 영향 점검에서 빼는 받는 쪽 열 개는 같은 값(동점) 가운데서 점(노드)의 순서대로 골랐어요. 영향 점검은 남은 트레이더에 대해 계약을 단위로 한 짝지은 부트스트랩이에요(다시 뽑기 2{,}000번, 시드 42).}
\label{tab:neutral-label-recipients}
\fitwidth{%
\begin{tabular}{lcc}
\toprule
 & 받는 쪽 & 다른 트레이더 \\
\midrule
\multicolumn{3}{l}{\emph{기간 W1, 2026년 7월--9월 (트레이더 $n=70$)}} \\
트레이더 수 & 37 & 33 \\
청산이 없는 비율 & 1.000 & 0.727 \\
청산 없이 닫은 비율의 평균 & 1.000 & 0.878 \\
이익을 낸 비율의 평균 & 0.325 & 0.479 \\
손실이 아닌 비율의 평균 & 0.433 & 0.644 \\
라벨 기간에 닫은 횟수의 중앙값 & 24.0 & 11.0 \\
\multicolumn{3}{l}{받은 허락 수가 가장 큰 받는 쪽 열 개를 뺀 P: $\Delta\tau=0.073$ [-0.039, 0.186]} \\
\midrule
\multicolumn{3}{l}{\emph{기간 W0, 2026년 4월--6월 (트레이더 $n=88$)}} \\
트레이더 수 & 40 & 48 \\
청산이 없는 비율 & 1.000 & 0.750 \\
청산 없이 닫은 비율의 평균 & 1.000 & 0.919 \\
이익을 낸 비율의 평균 & 0.362 & 0.474 \\
손실이 아닌 비율의 평균 & 0.402 & 0.607 \\
라벨 기간에 닫은 횟수의 중앙값 & 31.0 & 10.0 \\
\multicolumn{3}{l}{받은 허락 수가 가장 큰 받는 쪽 열 개를 뺀 P: $\Delta\tau=0.188$ [0.046, 0.331]} \\
\bottomrule
\end{tabular}
}
\end{table}


그냥 센 수의 수준에서 한 비교는 몇 안 되는 트레이더에 바탕을 두어요. W1에서 트레이더들 가운데 가장 큰 받은 허락 수는 1이에요. 그래서 P는 사실상 허락을 받은 트레이더 37개와 나머지 트레이더 33개를 견주는 거예요. W0에서는 가장 큰 값이 3이고, 트레이더 88개 가운데 40개가 허락을 받은 트레이더예요. 허락을 받은 트레이더는 두 기간 모두에서 하나도 빠짐없이 청산을 피했어요(37개 중 37개, 40개 중 40개). 나머지 트레이더는 W1에서 72.7\%, W0에서 75.0\%만 청산을 피했어요. 허락을 받은 트레이더는 포지션도 더 많이 닫았어요(중앙값이 W1에서 24 대 11, W0에서 31 대 10). 그리고 이익을 낸 비율과 손실이 아닌 비율이 더 낮았어요(W1에서 이익을 낸 비율의 평균 0.325 대 0.479). 이것이 R3 뒤에 있는 모양이에요. 받은 허락 수가 가장 큰, 허락을 받은 트레이더 10개를 빼면, P는 W1에서 $+0.073$ $[-0.039,0.186]$, W0에서 $+0.188$ $[0.046,0.331]$이에요. W1에서는 이 10개를 동점인 값들 가운데 점의 순서대로 골랐어요. 두 기간의 트레이더는 모두 계약 코드를 가지고 있어요. 그래서 계정 종류로는 허락을 받은 트레이더와 나머지를 가를 수 없어요. W1에서는 송금 수도 청산 없이 닫은 비율을 줄 세워요. 그리고 트레이더가 70개와 88개뿐이라서, 등록 문서가 주된 기간으로 정한 기간에서 허락한 주인 수가 송금 수보다 앞선다는 것은 가려지지 않아요.

\dissertationSubheading[sec:empirical-claims]{주장(C1--C6)}

이 장의 결과는 아래의 주장들을 뒷받침해요. 각 주장은 표~\ref{tab:hrq-claim-map}의 가설 열과 연구 질문 열에 이어져 있어요. 모든 주장은 수정 분석에 바탕을 둬요. C5의 규칙~A와~B의 판정, 그리고 C6의 R1--R3의 판정은 등록한 규칙을 고친 점수에 적용한 거예요. 이 적용은 등록 분석의 결과를 안 뒤에 했어요. 그래서 이 판정들은 재분석이에요. 등록 분석도 같은 판정에 이르렀어요(부록~\ref{sec:map-registered}).

\begin{enumerate}
\item[C1.] 비용. 같은 계약들의 전체 그래프에서, 풀이 프로그램의 시간은 그 그래프의 화살표 수를 따라가요. 부분 표본에서는 계약 10{,}000개 단계부터 그래요. 허락 그래프의 화살표 수는 송금 그래프의 44.8\%예요($T_{\mathrm{obs}}$에서 11,210,114개 대 25,001,216개). 그만큼 엔도스랭크를 한 번 푸는 데 드는 시간도 AWP보다 짧아요. AWP는 2.17배 오래 걸리고(점 7,882,356개에서 22.37~초 대 점 6,219,063개에서 10.29~초), 이때 화살표 수 비율은 2.23이에요(표~\ref{tab:benchmark-runtime}). C-PR(두 층을 합친 점 9,151,010개와 화살표 31,552,359개에서 36.16~초)과 S-PR(점 7,884,602개와 화살표 25,018,892개에서 34.42~초)은 송금 층을 걸어요. 그래서 AWP보다 오래 걸려요. 빠른 까닭은 알고리즘이 아니라, 더 듬성듬성한 그래프에 있어요.
\item[C2.] 같은 기간 일치. 각 점수는 자기가 걷는 화살표 종류의 대리 지표와 가장 잘 맞아요. 엔도스랭크는 허락 가족과 가장 잘 맞고($0.585$, $[0.578,0.592]$), AWP는 송금 가족과 가장 잘 맞아요($0.572$, $[0.564,0.580]$). 구간은 0에서 멀리 떨어져 있어요(표~\ref{tab:alignment-hybrid}). 이것은 의도한 구성개념 점검이에요. 그리고 그 일치의 일부는 만든 방식 때문에 저절로 생겨요. 금액을 달러로 재면, AWP는 들어온 송금의 금액 합을, 보낸 쪽의 수를 따라가는 것과 거의 같은 만큼 가깝게 따라가요(받은 양의 합 $0.566$, $[0.558,0.575]$; 입차수 $0.578$, $[0.569,0.587]$). 코호트의 모든 계약은 다섯 토큰의 허락을 받았어요. 그래서 엔도스랭크는 동점을 거의 남기지 않아요(서로 다른 점수 13,625개, 쌍의 0.17\%만 동점). 그리고 외톨이 점 규칙은 어떤 계수도 바꾸지 않아요(표~\ref{tab:tie-sensitivity}). C-PR의 섞은 점수들은 모든 가족에서 자기의 두 층 사이에 있어요.
\item[C3.] 갈라짐. 엔도스랭크와 AWP는 계약들을 서로 다르게 줄 세워요. $\tau_b=0.174$ $[0.163,0.185]$예요. 쌍의 96.7\%는 두 점수 모두에서 순서가 정해져 있어요. 그리고 그 가운데 41.2\%는 순서가 어긋나는 쌍이에요(표~\ref{tab:tie-sensitivity}, 그림~\ref{fig:rank-scatter}).
\item[C4.] 기간 밖에서의 차수 기준선. 두 걷기는 자기가 매끄럽게 다듬는 수와 견주었을 때 서로 다르게 움직여요. AWP는 두 기간 모두에서 나중의 새 송금자를 원시 송금 입차수보다 더 잘 줄 세워요(W0 $+0.043$, $[0.038,0.048]$; W1 $+0.022$, $[0.017,0.027]$). 그러니까 송금 쪽에서는, 시간 무게와 금액 무게, 그리고 활동 재시작을 쓰는 걷기가 그냥 센 수에 정보를 보탰어요. 계획에서 정한 고르게 다시 출발하기를 쓰는 엔도스랭크는, 두 기간 모두에서 새 허락 쌍을 원시 받은 허락 수보다 못하게 줄 세워요(W0 $-0.019$, $[-0.026,-0.011]$; W1 $-0.034$, $[-0.038,-0.029]$). AWP의 활동 재시작을 쓴 변형은 두 기간 모두에서 새 허락 쌍을 차수보다 더 잘 줄 세워요(W0 $+0.059$, $[0.050,0.067]$; W1 $+0.027$, $[0.021,0.034]$). 하지만 그 비교는 라벨을 본 뒤에 했어요. 그래서 이것은 발견이 아니라 가설이에요(\ref{sub:endorserank-restarts}절, 표~\ref{tab:fresh-posthoc}). 모든 차수 대비는 라벨을 본 뒤에 계산했어요. 계획에도 등록 문서에도 이름이 오른 대비는 하나도 없어요. W0에서 C-PR ($\lambda=1$)이 받은 허락 수보다 더 얻은 몫($-0.027$, $[-0.034,-0.019]$)과 S-PR의 두 대비도 라벨을 본 뒤에 계산했어요. W0을 송금을 받은 적이 있는 스펜더 14,094개로 좁혀도, 두 차수 비교는 그대로예요(엔도스랭크 $-0.013$, $[-0.020,-0.006]$; AWP $+0.056$, $[0.049,0.064]$; 표~\ref{tab:holdout-receiving}).
\item[C5.] 결합. 고친 점수에서, 두 종류의 화살표를 함께 걷는 걷기 하나는 한 층만 걷는 걷기들을 이기지 못해요. 규칙~A는 실패해요. C-PR은 새 허락 쌍에서 허락 걷기보다 나쁘지 않지만($+0.029$, $[0.022,0.036]$), 새 송금자에서는 송금 걷기보다 나빠요($-0.034$, $[-0.040,-0.029]$). 그리고 같은 기간 부분도 실패해요(표~\ref{tab:holdout-diff}과~\ref{tab:tau-diff-hybrid}). 규칙~B는 충족되지 않아요(표~\ref{tab:fresh-contrasts}). F1은 성립해요(새 허락 쌍에서 송금 걷기보다 $+0.110$, $[0.098,0.122]$). 그러니까 라벨을 본 뒤에 계산한 W0의 이득 $+0.092$($[0.080,0.105]$)가 다시 나타나요. F2는 실패해요($-0.028$, $[-0.033,-0.023]$로, 허용 폭 0.01을 넘어요). 그리고 F3은 아슬아슬하게 실패해요(트레이더의 청산 없이 닫은 비율에서 $+0.000$, $[-0.014,0.010]$이에요. 점 추정값은 둘 다 $0.350$으로 같지만, 아래 끝이 $-0.01$보다 아래에 있어요). 첫 번째 민감도 분석에서 AWP를 비교 대상으로 삼으면, F2는 뚜렷하게 실패해요($-0.067$, $[-0.072,-0.061]$). 라벨을 본 뒤에 계산한 W0의 대비 두 개에서, S-PR은 새 송금자에서 송금 걷기와 비슷해요($-0.001$, $[-0.004,0.002]$). 그리고 새 허락 쌍에서는 허락 걷기보다 아래에 있어요($-0.051$, $[-0.066,-0.038]$). 허락 화살표는 걷기를 실어 나를 때만 도움이 돼요. 혼합 점수들이 같은 기간 안에서 어디에 있는지는 \ref{sec:coupled-results}절에서 보고해요.
\item[C6.] 두 종류의 화살표 어느 쪽으로도 만들지 않은 결과 값. GMX~V2 계약 트레이더의 청산 없이 닫은 비율에서, 그냥 센 수의 수준에서 허락한 주인 수가 송금 수보다 앞선다는 것(R1)은 주된 기간 W1에서 뒷받침되지 않아요(P $=+0.084$, 본페로니 97.5\% 구간 $[-0.026,0.188]$; 층화하면 $+0.101$, $[-0.024,0.231]$). 다만 W0에서는 P가 양수예요($+0.182$, 97.5\% 구간 $[0.032,0.338]$). 페이지랭크가 앞선다는 것(R2)도 뒷받침되지 않아요(Q $=-0.041$, 97.5\% 구간 $[-0.148,0.064]$; W0에서 $-0.014$, 97.5\% 구간 $[-0.161,0.151]$). R3은 적용돼요. 허락한 주인 수는 이익을 낸 비율($-0.196$, $[-0.364,-0.046]$)과 손실이 아닌 비율($-0.252$, $[-0.391,-0.120]$)에서 트레이더를 송금 수보다 낮게 줄 세워요. 허락을 받은 적이 있는 트레이더는 두 기간 모두에서 하나도 빠짐없이 청산을 피했어요(37개 중 37개, 40개 중 40개). 나머지 트레이더는 72.7\%와 75.0\%만 청산을 피했어요. 하지만 표본이 작아요(계약 트레이더 70개와 88개). 그리고 송금 수도 이 라벨을 줄 세워요(W1에서 $0.310$, $[0.184,0.425]$)(표~\ref{tab:neutral-label-contrasts}과~\ref{tab:neutral-label-recipients}).
\end{enumerate}

순위 갈라짐은 RQ1에 답해요. 같은 기간 표들은 RQ2에, 벤치마크 표들은 RQ3에, 홀드아웃들은 RQ4와 RQ5에 답해요. C6은 연구 질문을 더하지 않아요. 엔도스랭크와 AWP를 직접 견준 비교는 부록~\ref{ch:appendix-er-awp}와 표~\ref{tab:neutral-label-contrasts}에 있어요. 그 가운데 어느 것도 이 논문의 주장을 담고 있지 않아요. \ref{ch:discussion}장에서는 이 발견들을 주된 연구 질문에 비추어 읽어요. 주된 연구 질문은, 아비트럼 원의 스마트 계약에 매기는 평판 점수로서 ERC-20 허락 화살표가 페이지랭크에 무엇을 보태는지예요. 이때 송금 걷기와 견주고, 두 걷기가 매끄럽게 다듬는 원시 차수와도 견줘요. 그 장에서는 이 발견들을 \ref{ch:introduction}장의 목표에도 비추어 읽어요.
